\documentclass[11pt]{article}
\usepackage{fontspec}
% Charter et Menlo (macOS) ; ailleurs, XCharter et DejaVu Sans Mono, dont ils dérivent.
\IfFontExistsTF{Charter}{\setmainfont{Charter}}{\setmainfont{XCharter}}
\IfFontExistsTF{Menlo}{\setmonofont{Menlo}[Scale=0.82]}{\setmonofont{DejaVu Sans Mono}[Scale=0.82]}
\usepackage{polyglossia}\setmainlanguage{french}\setotherlanguage{english}
\usepackage{microtype}
% L’ : la paire de crénage de Charter colle l’apostrophe au L ; on rajoute une fine espace.
\XeTeXinterchartokenstate=1
\newXeTeXintercharclass\LetterL
\newXeTeXintercharclass\Apostrophe
\XeTeXcharclass`L=\LetterL
\XeTeXcharclass`'=\Apostrophe
\XeTeXcharclass"2019=\Apostrophe
\XeTeXinterchartoks\LetterL\Apostrophe={\kern0.07em}
\usepackage[a4paper,margin=27mm]{geometry}
\usepackage{setspace}\setstretch{1.1}
\usepackage{xcolor}
\usepackage[hidelinks]{hyperref}
\usepackage{amsmath,amssymb}
\definecolor{frise}{HTML}{2F4DA8}

\newcommand{\q}[1]{«\,#1\,»}
% La vidéo de la section en cours ; chaque entrée renvoie à sa seconde.
\newcommand{\videoid}{}
\newcommand{\video}[1]{\renewcommand{\videoid}{#1}}
\newcommand{\face}[1]{\par\medskip\noindent\textbf{#1}\par\smallskip}
% Une entrée : le moment de la bande (lien vers YouTube), puis l'énoncé.
\newlength{\horo}\settowidth{\horo}{\small\texttt{00:00:00}\quad}
\newcommand{\entree}[3]{\par\noindent\hangindent=\horo\hangafter=1%
  \makebox[\horo][l]{\small\href{https://www.youtube.com/watch?v=\videoid&t=#1s}{\textcolor{frise}{\texttt{#2}}}}#3\par\smallskip}
\raggedbottom\emergencystretch=1.5em

\title{\Large Cours de Buffalo sur les topos (1973) :\\ index des résultats}
\author{Michel Hua\\ \small Chercheur indépendant, Paris\\ \small \texttt{michel@commutator.io}}
\date{\small Novembre 2026}

\begin{document}
\maketitle

\section*{Introduction}\addcontentsline{toc}{section}{Introduction}

\emph{Le cours.} D'avril à juillet 1973, Alexandre Grothendieck a fait à l'université d'État de New York à Buffalo un cours sur la théorie des topos, en anglais, à côté de deux autres cours, sur la géométrie algébrique et sur les groupes algébriques. Les séances ont été enregistrées sur bandes magnétiques. Ce document s'appuie sur les enregistrements publiés par la chaîne YouTube \q{Erwin Schulhoff} en quatre vidéos ; il cite chaque séance par sa date, sa bande et sa face, d'après les descriptions de ces vidéos, et chaque moment par sa position dans la vidéo.

\emph{Ce qu'est cet index.} Pour chaque résultat, définition ou exemple traité dans une séance, l'index donne une ligne : l'énoncé, dans la langue mathématique d'aujourd'hui et dans nos propres mots, et le moment de la bande où il est traité, sous la forme d'un lien vers la vidéo à la seconde près. Une entrée correspond à un énoncé, non à une minute de cours. Quelques mots entre guillemets reprennent un nom ou un terme forgé dans le cours, quand il importe. Quand ce qui est dit est esquissé ou incomplet, l'entrée donne l'énoncé correct, avec les hypothèses qu'il demande ; une affirmation fausse telle qu'elle est entendue n'est pas reprise. Un passage dont l'audition reste douteuse n'est pas indexé.

\emph{Ce qu'il n'est pas.} Cet index n'est pas une transcription ni un résumé du cours. Il ne reprend ni les explications, ni les exemples développés, ni les digressions, ni l'ordre de l'exposé au-delà de la succession des énoncés, et il ne cite pas le cours. Le cours tel qu'il a été prononcé est une œuvre orale protégée, dont les droits appartiennent aux ayants droit d'Alexandre Grothendieck ; seuls les enregistrements font foi, et l'on s'y reporte par les liens. Les énoncés sont ici des énoncés mathématiques standard, rédigés par nous.

\emph{Ce qui manque.} Les cinq séances de la deuxième vidéo (30 mai, 4, 6, 20 et 27 juin 1973) ne sont pas encore indexées : elles sont marquées \emph{à venir}. Quatre faces au son trop dégradé ne le sont pas non plus (4 avril, bande 1, face B ; 11 mai, bande 2, face A ; et deux faces sans étiquette ni date à la fin de la quatrième vidéo).

\tableofcontents

\section{4 avril 1973}\label{s-1973-04-04}
\video{URU5jgTvjhU}
\noindent\textit{Vidéo 1/4 (\href{https://www.youtube.com/watch?v=URU5jgTvjhU}{\texttt{URU5jgTvjhU}}) : bande 1, face A (0:00:00--0:49:11); bande 1, face B (0:49:11--1:36:58); bande 2, face A (1:36:58--2:28:43); bande 2, face B (2:28:43--2:33:29).}\par\medskip
\face{Bande 1, face A}
\entree{206}{0:03:26}{Les ouverts $\mathcal{O}(X)$ d'un espace $X$, ordonnés par inclusion, forment une catégorie : $\mathrm{Hom}(U,V)$ est un singleton si $U\subset V$, vide sinon.}
\entree{335}{0:05:35}{Définition : un préfaisceau d'ensembles sur $X$ est un foncteur $F\colon\mathcal{O}(X)^{\mathrm{op}}\to\mathbf{Ens}$ ; les applications $F(V)\to F(U)$, $U\subset V$, sont les restrictions.}
\entree{466}{0:07:46}{Un morphisme de préfaisceaux $F\to G$ est une transformation naturelle : des applications $F(U)\to G(U)$ compatibles aux restrictions.}
\entree{596}{0:09:56}{Définition : $F$ est séparé si, pour tout recouvrement ouvert $U=\bigcup U_i$, l'application $F(U)\to\prod_i F(U_i)$ est injective.}
\entree{859}{0:14:19}{Définition : $F$ est un faisceau si, pour tout recouvrement $U=\bigcup U_i$, $F(U)\to\prod_i F(U_i)\rightrightarrows\prod_{i,j}F(U_i\cap U_j)$ est exact (noyau).}
\entree{988}{0:16:28}{Les faisceaux sur $X$ forment une sous-catégorie pleine $\mathrm{Sh}(X)$ (notée ici \q{Top(X)}) de la catégorie des préfaisceaux.}
\entree{1127}{0:18:47}{Exemple : pour un ensemble $E$, le préfaisceau $U\mapsto E^U$ des applications de $U$ dans $E$ est un faisceau.}
\entree{1254}{0:20:54}{Définition : un sous-préfaisceau $G\subset F$ est une famille de parties $G(U)\subset F(U)$ stables par les restrictions.}
\entree{1452}{0:24:12}{Un sous-préfaisceau d'un préfaisceau séparé est séparé.}
\entree{1452}{0:24:12}{Un sous-préfaisceau $G$ d'un faisceau $F$ est un faisceau si et seulement si l'appartenance d'une section à $G$ est une propriété locale sur $U$.}
\entree{1579}{0:26:19}{Exemple : si $E$ est un espace topologique, les applications continues des ouverts de $X$ dans $E$ forment un sous-faisceau de $U\mapsto E^U$.}
\entree{1643}{0:27:23}{Exemple : pour des variétés différentiables (ou analytiques) $X$ et $E$, les applications différentiables (resp. analytiques) forment un faisceau.}
\entree{1769}{0:29:29}{Exemple : pour $p\colon E\to X$ continue, les sections continues $\Gamma(U,E)$ au-dessus des ouverts $U$ forment un faisceau.}
\entree{1965}{0:32:45}{Le faisceau des applications continues de $X$ dans $E$ est isomorphe au faisceau des sections de la seconde projection $E\times X\to X$.}
\entree{2099}{0:34:59}{Exemple : sur un espace localement compact $X$, les mesures de Radon sur les ouverts de $X$ forment un faisceau.}
\entree{2164}{0:36:04}{Exemple : sur une variété différentiable, les distributions sur les ouverts forment un faisceau.}
\entree{2164}{0:36:04}{Exemple : $U\mapsto\{\text{parties de }U\}$, restriction par intersection, est un faisceau ; les parties fermées dans $U$ en forment un sous-faisceau.}
\entree{2291}{0:38:11}{Exemple : sur un espace analytique, les sous-ensembles analytiques des ouverts forment un faisceau.}
\entree{2487}{0:41:27}{Le faisceau des sections $\widetilde{E}$ d'un espace $E$ au-dessus de $X$ définit un foncteur $\mathrm{Fib}(X)\to\mathrm{Sh}(X)$, $f\mapsto(s\mapsto f\circ s)$.}
\entree{2680}{0:44:40}{Ce foncteur n'est pas pleinement fidèle : pour $X$ un point, il s'identifie au foncteur d'oubli des espaces topologiques vers $\mathbf{Ens}$.}
\entree{2745}{0:45:45}{Sur un point, le foncteur $\Gamma\colon F\mapsto F(\mathrm{pt})$ est une équivalence de catégories $\mathrm{Sh}(\mathrm{pt})\simeq\mathbf{Ens}$.}
\entree{2811}{0:46:51}{Pour $U$ ouvert, $\Gamma_U\colon F\mapsto F(U)$ est un foncteur $\mathrm{Sh}(X)\to\mathbf{Ens}$ ; on note $\Gamma(F)=\Gamma(X,F)=F(X)$.}
\face{Bande 1, face B}
\noindent\textit{Non indexée (son trop dégradé pour être transcrit).}\par
\face{Bande 2, face A}
\entree{5826}{1:37:06}{Les notions duales des produits, noyaux de doubles flèches, produits fibrés et objet final sont les sommes, conoyaux, sommes amalgamées et objet initial.}
\entree{5888}{1:38:08}{Dans $\mathbf{Ens}$, les sommes sont les réunions disjointes et l'objet initial est $\varnothing$.}
\entree{5888}{1:38:08}{Dans $\mathbf{Ens}$, le conoyau de $E\rightrightarrows F$ est le quotient de $F$ par la relation d'équivalence engendrée par l'image de $E$ dans $F\times F$.}
\entree{6015}{1:40:15}{Exemple : dans les groupes, la somme est le produit libre ; la somme de $n$ copies de $\mathbb{Z}$ est le groupe libre à $n$ générateurs.}
\entree{6082}{1:41:22}{Pour les structures algébriques usuelles, les limites projectives se calculent sur les ensembles sous-jacents, mais pas les limites inductives.}
\entree{6212}{1:43:32}{Si $S$ admet les limites (resp. colimites) d'un type donné, $\mathrm{Hom}(C^{\mathrm{op}},S)$ les admet aussi, calculées argument par argument.}
\entree{6275}{1:44:35}{La catégorie opposée à $\mathrm{Hom}(C^{\mathrm{op}},S)$ est $\mathrm{Hom}(C,S^{\mathrm{op}})$.}
\entree{6401}{1:46:41}{Les préfaisceaux sur $X$ admettent toutes les limites inductives, calculées ouvert par ouvert.}
\entree{6469}{1:47:49}{La sous-catégorie $\mathrm{Sh}(X)$ n'est stable dans les préfaisceaux pour aucun type usuel de limite inductive.}
\entree{6602}{1:50:02}{Exemple : la somme, prise dans les préfaisceaux, de l'objet final avec lui-même (préfaisceau constant à deux éléments) n'est pas un faisceau.}
\entree{6730}{1:52:10}{Pour tout faisceau $F$, $F(\varnothing)$ est un singleton ($\varnothing$ étant recouvert par la famille vide).}
\entree{6797}{1:53:17}{Exemple : le préfaisceau initial $U\mapsto\varnothing$ n'est pas un faisceau.}
\entree{6861}{1:54:21}{Une limite inductive de faisceaux existe si et seulement si le foncteur $G\mapsto\mathrm{Hom}(P,G)$ sur $\mathrm{Sh}(X)$ est représentable, $P$ la limite prise dans les préfaisceaux.}
\entree{7053}{1:57:33}{Définition : le faisceau associé à un préfaisceau $P$ est un faisceau $\widetilde{P}$ muni de $P\to\widetilde{P}$ universel pour les morphismes de $P$ vers des faisceaux.}
\entree{7116}{1:58:36}{Théorème : tout préfaisceau admet un faisceau associé, unique à isomorphisme unique près.}
\entree{7310}{2:01:50}{Corollaire : $\mathrm{Sh}(X)$ admet toutes les limites inductives, faisceau associé à la limite inductive calculée dans les préfaisceaux.}
\entree{7310}{2:01:50}{Un préfaisceau et son faisceau associé ont les mêmes fibres.}
\entree{7450}{2:04:10}{En termes d'espaces étalés, la somme d'une famille de faisceaux est la somme topologique disjointe, qui est étalée sur $X$.}
\entree{7516}{2:05:16}{Le conoyau de $F\rightrightarrows G$ (espaces étalés sur $X$) est $G/R$, muni de la topologie quotient, $R$ la relation engendrée ; $G/R$ est étalé sur $X$.}
\entree{7642}{2:07:22}{Les sommes et les conoyaux de doubles flèches suffisent à construire toutes les limites inductives.}
\entree{7704}{2:08:24}{Principe : $\mathrm{Sh}(X)$ a les propriétés d'exactitude de $\mathbf{Ens}$ qui font intervenir des limites inductives quelconques et des limites projectives finies.}
\entree{7768}{2:09:28}{Dans $\mathbf{Ens}$, un produit quelconque d'épimorphismes est un épimorphisme ; c'est faux dans $\mathrm{Sh}(X)$ en général.}
\entree{7960}{2:12:40}{Les catégories de préfaisceaux (de foncteurs vers $\mathbf{Ens}$) ont toutes les propriétés d'exactitude de $\mathbf{Ens}$, limites et colimites étant calculées argument par argument.}
\entree{8025}{2:13:45}{Le foncteur $a\colon P\mapsto\widetilde{P}$ est adjoint à gauche de l'inclusion $i$ : $\mathrm{Hom}(P,iF)\simeq\mathrm{Hom}(aP,F)$ ; il commute aux limites inductives.}
\entree{8089}{2:14:49}{Théorème : le foncteur faisceau associé est exact à gauche (il commute aux produits finis et aux noyaux de doubles flèches).}
\entree{8288}{2:18:08}{Dans $\mathrm{Sh}(X)$, les conoyaux sont stables par changement de base : $\mathrm{Coker}(F\rightrightarrows G)\times_H H'$ est le conoyau des images réciproques sur $H'$.}
\entree{8672}{2:24:32}{Dans $\mathrm{Sh}(X)$, les sommes sont stables par changement de base : $(\coprod_i F_i)\times_F F'\simeq\coprod_i(F_i\times_F F')$.}
\entree{8863}{2:27:43}{Définition : un objet quotient de $F$ est une classe d'épimorphismes $F\to G$ à isomorphisme près sous $F$.}
\face{Bande 2, face B}
\entree{8923}{2:28:43}{Pour tout morphisme $F\to G$, le produit fibré $F\times_G F\subset F\times F$ est une relation d'équivalence sur $F$.}
\entree{8987}{2:29:47}{À une relation d'équivalence $R\subset F\times F$ est associé le conoyau de $R\rightrightarrows F$, quotient de $F$ (épimorphisme).}
\entree{9052}{2:30:52}{Définition : un quotient est effectif s'il est le quotient de sa relation $F\times_G F$ ; une relation $R$ est effective si $R=F\times_{F/R}F$.}
\entree{9117}{2:31:57}{Dans $\mathrm{Sh}(X)$, comme dans $\mathbf{Ens}$, relations d'équivalence et quotients sont effectifs et se correspondent bijectivement.}
\section{11 mai 1973}\label{s-1973-05-11}
\video{URU5jgTvjhU}
\noindent\textit{Vidéo 1/4 (\href{https://www.youtube.com/watch?v=URU5jgTvjhU}{\texttt{URU5jgTvjhU}}) : bande 1, face A (2:33:29--3:22:55); bande 1, face B (3:22:55--4:02:08); bande 2, face A (4:02:08--4:48:29); bande 2, face B (4:48:29--5:36:41).}\par\medskip
\face{Bande 1, face A}
\entree{9213}{2:33:33}{Dans $\mathbf{Ens}$, $G$ opère simplement transitivement sur $E$ si $G\times E\to E\times E$, $(g,x)\mapsto(gx,x)$, est bijective ; $E$ est un torseur si de plus $E\to\mathrm{pt}$ est surjective.}
\entree{9274}{2:34:34}{$E\to F$ est un monomorphisme si et seulement si $E\to E\times_F E$ est un isomorphisme ; dualement, un épimorphisme si $F\sqcup_E F\to F$ est un isomorphisme.}
\entree{9338}{2:35:38}{Définition : un torseur sous un faisceau de groupes $G$ est un $G$-faisceau $E$ avec $G\times E\simeq E\times E$ et $E\to 1$ épimorphisme (et non $E(U)$ torseur pour tout $U$).}
\entree{9402}{2:36:42}{Un morphisme de faisceaux est un épimorphisme si et seulement s'il l'est sur les fibres ; $E$ est un $G$-torseur si et seulement si chaque $E_x$ est un $G_x$-torseur.}
\entree{9530}{2:38:50}{Un torseur est trivial (isomorphe à $G$ opérant par translations) si et seulement s'il admet une section globale ; il en existe de non triviaux.}
\entree{9658}{2:40:58}{Principe : énoncés sur les structures définies par limites finies et colimites se vérifient fibre à fibre, les foncteurs fibres formant une famille conservative exacte.}
\entree{9900}{2:45:00}{Les structures de corps et d'anneau intègre ne s'expriment pas par des limites projectives finies seules.}
\entree{10229}{2:50:29}{Un faisceau d'anneaux $K$ est dit intègre si $K'\sqcup 1\to K$ (inclusion et section nulle) est un épimorphisme, $K'$ défini ci-dessous.}
\entree{10358}{2:52:38}{Pour deux faisceaux $K$, $L$, le préfaisceau $\underline{\mathrm{Hom}}(K,L)\colon U\mapsto\mathrm{Hom}(K|_U,L|_U)$ est un faisceau.}
\entree{10423}{2:53:43}{Les sections $f$ de $K$ dont les multiplications à gauche et à droite sont des monomorphismes de $K|_U$ forment un sous-faisceau $K'$ (propriété locale).}
\entree{10756}{2:59:16}{Définition : $(Z_i)$ est une famille génératrice de $C$ si tout monomorphisme $E\to F$ tel que $\mathrm{Hom}(Z_i,E)\to\mathrm{Hom}(Z_i,F)$ soit bijectif pour tout $i$ est un isomorphisme.}
\entree{10884}{3:01:24}{Dans $\mathrm{Sh}(X)$, $(Z_i)$ est génératrice si et seulement si la famille de foncteurs $\mathrm{Hom}(Z_i,-)$ est conservative.}
\entree{11028}{3:03:48}{Dans $\mathrm{Sh}(X)$, tout morphisme se factorise de façon essentiellement unique en un épimorphisme suivi d'un monomorphisme (image).}
\entree{11157}{3:05:57}{Dans $\mathrm{Sh}(X)$, un morphisme à la fois monomorphisme et épimorphisme est un isomorphisme.}
\entree{11349}{3:09:09}{Un ouvert $U$ de $X$ définit un faisceau : $U(V)$ est un singleton si $V\subset U$, vide sinon.}
\entree{11484}{3:11:24}{Pour tout faisceau $E$, $\mathrm{Hom}(U,E)=E(U)$ ; les ouverts de $X$ forment donc une petite famille génératrice de $\mathrm{Sh}(X)$.}
\entree{11555}{3:12:35}{Les faisceaux abéliens sur $X$ (objets en groupes abéliens de $\mathrm{Sh}(X)$) forment une catégorie abélienne $\mathrm{Ab}(X)$.}
\entree{11687}{3:14:47}{$\mathrm{Ab}(X)$ admet limites projectives et inductives ; les limites inductives filtrantes y commutent à l'oubli et sont exactes.}
\entree{11816}{3:16:56}{Définition : un faisceau abélien $I$ est injectif si tout morphisme $F\to I$ se prolonge le long de tout monomorphisme $F\to G$.}
\entree{11881}{3:18:01}{Une catégorie abélienne à limites inductives filtrantes exactes et à petite famille génératrice a assez d'injectifs ; c'est le cas de $\mathrm{Ab}(X)$ (faisceaux abéliens libres sur les $Z_i$).}
\entree{11945}{3:19:05}{Pour un faisceau d'anneaux $A$, $\mathrm{Mod}(A)$ est abélienne, à colimites filtrantes exactes et générateurs, donc a assez d'injectifs ($\mathrm{Ext}$, cohomologie).}
\entree{12072}{3:21:12}{En général $\mathrm{Ab}(X)$ n'a pas assez de projectifs ; les limites projectives filtrantes n'y sont pas exactes.}
\face{Bande 1, face B}
\entree{12248}{3:24:08}{Définition : pour $f\colon X\to Y$, un $f$-morphisme $G\to F$ est une famille $G(V)\to F(U)$, pour $f(U)\subset V$, compatible aux restrictions.}
\entree{12440}{3:27:20}{Un $f$-morphisme équivaut à une famille $G(V)\to F(f^{-1}(V))$ compatible aux restrictions.}
\entree{12633}{3:30:33}{$V\mapsto f^{-1}(V)$ est un foncteur croissant $\mathcal{O}_Y\to\mathcal{O}_X$ qui commute aux réunions quelconques.}
\entree{12702}{3:31:42}{Définition : image directe d'un préfaisceau $F$ sur $X$, $f_*F(V)=F(f^{-1}(V))$, préfaisceau sur $Y$.}
\entree{12768}{3:32:48}{Les $f$-morphismes $G\to F$ correspondent bijectivement aux morphismes $G\to f_*F$ de préfaisceaux sur $Y$.}
\entree{12832}{3:33:52}{Sur les préfaisceaux, $f_*$ commute aux limites projectives et inductives quelconques.}
\entree{12899}{3:34:59}{L'image directe d'un faisceau est un faisceau.}
\entree{13097}{3:38:17}{$f_*\colon\mathrm{Sh}(X)\to\mathrm{Sh}(Y)$ commute aux limites projectives quelconques.}
\entree{13162}{3:39:22}{En général $f_*$ ne commute ni aux sommes ni aux épimorphismes : pour $Y$ un point, $f_*$ est $\Gamma(X,-)$.}
\entree{13294}{3:41:34}{Définition : l'image inverse $f^*G$ est munie d'un $f$-morphisme universel $G\to f^*G$ ; $f^*$ est l'adjoint à gauche de $f_*$.}
\entree{13485}{3:44:45}{Pour tout foncteur $\varphi\colon C\to D$, la restriction $\widehat{D}\to\widehat{C}$, $P\mapsto P\circ\varphi$, a un adjoint à gauche et un adjoint à droite.}
\entree{13614}{3:46:54}{$f^*G$ est le faisceau associé à l'image inverse préfaisceau $f^*_{\mathrm{pr}}G$.}
\entree{13743}{3:49:03}{$f^*$ commute aux limites inductives quelconques (adjoint à gauche).}
\entree{13813}{3:50:13}{Théorème : $f^*$ est exact à gauche (il commute aux limites projectives finies).}
\entree{13945}{3:52:25}{L'image réciproque d'un espace $G\to Y$ par $f$ est $X\times_Y G\subset X\times G$, topologie induite ; c'est un produit fibré dans les espaces topologiques.}
\entree{14011}{3:53:31}{La propriété d'être étalé est stable par changement de base.}
\entree{14075}{3:54:35}{Pour les faisceaux vus comme espaces étalés, $f^*G$ correspond au produit fibré $X\times_Y G$.}
\entree{14268}{3:57:48}{Les espaces étalés sur $Y$ sont stables par limites projectives finies dans les espaces au-dessus de $Y$.}
\entree{14268}{3:57:48}{La fibre de $X\times_Y G$ en $x$ est homéomorphe à la fibre de $G$ en $f(x)$.}
\entree{14332}{3:58:52}{$(f^*G)_x\simeq G_{f(x)}$ ; d'où à nouveau l'exactitude à gauche de $f^*$.}
\face{Bande 2, face A}
\noindent\textit{Non indexée (son trop dégradé pour être transcrit).}\par
\face{Bande 2, face B}
\entree{17311}{4:48:31}{Deux morphismes $u,v\colon F\to G$ de faisceaux sont égaux si et seulement si $u_x=v_x$ pour tout $x$ : les foncteurs fibres forment une famille fidèle.}
\entree{17440}{4:50:40}{Les foncteurs fibres commutent aux limites inductives et aux limites projectives finies et forment une famille conservative.}
\entree{17585}{4:53:05}{$F\to G$ est un épimorphisme de faisceaux si et seulement si toute section de $G$ se relève localement en une section de $F$.}
\entree{17712}{4:55:12}{Un épimorphisme de faisceaux n'est pas un épimorphisme de préfaisceaux en général : $\Gamma(F)\to\Gamma(G)$ n'est pas toujours surjectif.}
\entree{17712}{4:55:12}{Pour $X$ compact, $\Gamma(X,-)$ transforme les épimorphismes de faisceaux en surjections si et seulement si $X$ est totalement discontinu.}
\entree{17970}{4:59:30}{Définition : une loi de composition sur un objet $E$ de $C$ est une flèche $E\times E\to E$ ; commutativité et associativité s'expriment par des diagrammes.}
\entree{18102}{5:01:42}{Une unité à gauche est une flèche $1\to E$ ($1$ final) telle que $E\simeq 1\times E\to E\times E\to E$ soit l'identité.}
\entree{18295}{5:04:55}{Par le plongement de Yoneda, une loi $E\times E\to E$ équivaut à des lois sur les $E(Z)=\mathrm{Hom}(Z,E)$, fonctorielles en $Z$.}
\entree{18427}{5:07:07}{Définition : une structure algébrique (groupe, anneau, module, etc.) sur $E$ est une telle structure sur chaque $E(Z)$, fonctorielle en $Z$.}
\entree{18687}{5:11:27}{La structure de catégorie ($D_1\rightrightarrows D_0$, composition $D_1\times_{D_0}D_1\to D_1$, unités) se décrit par des produits fibrés.}
\entree{18814}{5:13:34}{Définition : une catégorie interne à $C$ est un couple $(D_0,D_1)$ et une structure de catégorie sur $(D_0(Z),D_1(Z))$, fonctorielle en $Z$.}
\entree{18881}{5:14:41}{Le plongement $C\to\widehat{C}$ commute aux limites projectives quelconques, mais pas aux limites inductives en général.}
\entree{18945}{5:15:45}{Une structure définie par limites projectives sur un faisceau $F$ équivaut à la même structure sur chaque $F(U)$, fonctorielle en $U$.}
\entree{19073}{5:17:53}{Un faisceau de groupes est un foncteur $\mathcal{O}(X)^{\mathrm{op}}\to\mathbf{Grp}$ dont le composé avec l'oubli est un faisceau.}
\entree{19341}{5:22:21}{Les axiomes d'une structure définie par limites projectives finies se vérifient fibre à fibre.}
\entree{19406}{5:23:26}{Un foncteur qui commute aux produits finis transforme objets en groupes (anneaux, etc.) en objets de même type.}
\entree{19532}{5:25:32}{Les fibres d'un faisceau de groupes (d'anneaux, etc.) sont des groupes (des anneaux, etc.).}
\entree{19665}{5:27:45}{Définition : pour une $k$-cogèbre $A$, un $A$-comodule est un $k$-module $M$ muni de $M\to M\otimes_k A$ satisfaisant coassociativité et counité.}
\entree{19856}{5:30:56}{Le produit tensoriel de modules se construit par des limites inductives (sommes, quotients), non par des limites projectives seules.}
\entree{19999}{5:33:19}{Pour des faisceaux de modules, $M(U)\otimes_{k(U)}A(U)\to(M\otimes_k A)(U)$ n'est en général ni injective ni surjective.}
\entree{20064}{5:34:24}{Définition : $E$ est un torseur sous un groupe $G$ si $E\neq\varnothing$ et $G$ opère simplement transitivement sur $E$.}
\section{18 mai 1973}\label{s-1973-05-18}
\video{URU5jgTvjhU}
\noindent\textit{Vidéo 1/4 (\href{https://www.youtube.com/watch?v=URU5jgTvjhU}{\texttt{URU5jgTvjhU}}) : bande 1, face A (5:36:41--6:26:00); bande 1, face B (6:26:00--7:14:32); bande 2, face A (7:14:32--8:03:53); bande 2, face B (8:03:53--8:23:34).}\par\medskip
\face{Bande 1, face A}
\entree{20360}{5:39:20}{Pour un point $y$ de $Y$ et l'inclusion $i_y:\{y\}\to Y$, on a $i_y^*F\simeq F_y$ : la fibre est l'image inverse par un point.}
\entree{20491}{5:41:31}{Si $f:X\to Y$ est étale, $X$ est un objet de $\mathrm{Sh}(Y)$ et $\mathrm{Sh}(X)\simeq\mathrm{Sh}(Y)/X$ ; un composé de morphismes étales est étale.}
\entree{20620}{5:43:40}{Définition : pour un objet $X$ d'une catégorie $C$, la catégorie $C/X$ a pour objets les flèches $X'\to X$ et pour morphismes les triangles commutatifs.}
\entree{20684}{5:44:44}{Si $f:X\to Y$ est étale, $f^*G$ est le produit $G\times X$ dans $\mathrm{Sh}(Y)$, vu comme objet de $\mathrm{Sh}(Y)/X$.}
\entree{20820}{5:47:00}{Pour $g:X'\to X$ dans $\mathrm{Sh}(Y)$, le foncteur $g^*:\mathrm{Sh}(Y)/X\to\mathrm{Sh}(Y)/X'$ est le changement de base $F\mapsto F\times_X X'$.}
\entree{21075}{5:51:15}{$f_*$ transforme toute structure définie par limites projectives (groupe, anneau, etc.) en structure de même espèce, mais pas les torseurs.}
\entree{21139}{5:52:19}{$f^*$ respecte les structures définies par limites projectives finies et limites inductives : l'image inverse d'un $G$-torseur est un $f^*G$-torseur.}
\entree{21393}{5:56:33}{Définition : un objet anneau $A$ de $C$ est un objet muni d'une structure d'anneau sur chaque $A(Z)=\mathrm{Hom}_C(Z,A)$, fonctorielle en $Z$.}
\entree{21458}{5:57:38}{Les inversibles $A^*(Z)\subset A(Z)$ forment un foncteur en groupes, représenté par le produit fibré de $m:A\times A\to A$ et de $1:e\to A$.}
\entree{21587}{5:59:47}{Définition : un objet anneau $A$ est un objet corps si le morphisme $A^*\sqcup e\to A$ (inclusion et section nulle) est un isomorphisme.}
\entree{21778}{6:02:58}{Intégrité, première condition : pour $C=m^{-1}(0)\subset A\times A$, le morphisme $A\sqcup A\to C$, $x\mapsto(x,0)$, $x\mapsto(0,x)$, est épimorphique.}
\entree{22035}{6:07:15}{Un morphisme $S\to C$ est un épimorphisme ssi la flèche canonique $C\sqcup_S C\to C$ est un isomorphisme.}
\entree{22099}{6:08:19}{Intégrité, seconde condition ($0\neq1$) : le noyau $e'\subset e$ de $0,1:e\rightrightarrows A$ est initial ; dans $\mathrm{Sh}(X)$, $e'$ est le plus grand ouvert où $A=0$.}
\entree{22292}{6:11:32}{Toute espèce de structure définie par limites projectives finies et limites inductives quelconques admet un \q{topos classifiant}, muni d'un modèle universel.}
\entree{22357}{6:12:37}{Les anneaux normaux, les anneaux réduits, les anneaux de spectre connexe se définissent aussi par de tels axiomes diagrammatiques.}
\entree{22420}{6:13:40}{Un faisceau d'anneaux n'est jamais un préfaisceau de corps : sa valeur sur $\varnothing$ est l'anneau nul.}
\entree{22545}{6:15:45}{Sur un espace, un tel axiome se vérifie fibre à fibre : un faisceau d'anneaux $A$ est un objet corps ssi chaque fibre $A_x$ est un corps.}
\entree{22609}{6:16:49}{Le faisceau constant $I_X=p^*I$ ($p:X\to\mathrm{pt}$) a pour espace étalé $I\times X$ ; c'est la somme de $I$ copies de l'objet final.}
\entree{22740}{6:19:00}{Si $I$ est un groupe, un anneau ou un corps, $I_X$ est un faisceau de groupes, d'anneaux ou de corps ; par exemple $K_X$ pour un corps $K$.}
\entree{22805}{6:20:05}{Sur un espace analytique complexe réduit et localement irréductible, le faisceau $\mathcal{M}_X$ des fonctions méromorphes est un faisceau de corps.}
\entree{22867}{6:21:07}{La fibre $\mathcal{M}_{X,x}$ est l'anneau total des fractions de $\mathcal{O}_{X,x}$ ; c'est un corps dès que $\mathcal{O}_{X,x}$ est intègre.}
\entree{22934}{6:22:14}{Sur un schéma intègre, le faisceau des fonctions rationnelles est un faisceau de corps constant ; le faisceau $\mathcal{M}_X$ analytique ne l'est pas.}
\entree{22998}{6:23:18}{Définition : $\mathcal{H}om(F,G)(U)=\mathrm{Hom}(F|_U,G|_U)$ définit un faisceau sur $X$, à distinguer de l'ensemble $\mathrm{Hom}(F,G)$.}
\face{Bande 1, face B}
\entree{23169}{6:26:09}{Pour $X'$ étale sur $X$ : $\mathrm{Hom}(X',\mathcal{H}om(F,G))\simeq\mathrm{Hom}_{X'}(F\times X',G\times X')\simeq\mathrm{Hom}(F\times X',G)$.}
\entree{23366}{6:29:26}{Définition : dans une catégorie à produits, $\mathcal{H}om(F,G)$ est, s'il existe, l'objet représentant $X'\mapsto\mathrm{Hom}(F\times X',G)$ ; il existe dans $\mathrm{Sh}(X)$.}
\entree{23428}{6:30:28}{Pour des objets groupes $F,G$, $\mathcal{H}om_{\mathrm{gr}}(F,G)$ représente $X'\mapsto\mathrm{Hom}_{\mathrm{gr}}(F_{X'},G_{X'})$, homomorphismes d'objets groupes de $C/X'$.}
\entree{23555}{6:32:35}{Le faisceau $\mathcal{P}(F)$ : $\mathcal{P}(F)(U)$ est l'ensemble des sous-faisceaux de $F|_U$ ; pour $X$ réduit à un point, c'est l'ensemble des parties de $F$.}
\entree{23684}{6:34:44}{$\mathcal{P}(F)$ représente $X'\mapsto\mathrm{Sub}(F\times X')$ ; les sous-objets d'un objet de $C/X'$ sont les mêmes dans $C$ et dans $C/X'$.}
\entree{23881}{6:38:01}{Pour l'objet final $e$, $\mathcal{P}(e)$ est le faisceau qui associe à $U$ l'ensemble de ses ouverts ; il représente $X'\mapsto\mathrm{Sub}(X')$.}
\entree{24070}{6:41:10}{Si $\mathcal{P}(e)$ existe, $\mathcal{P}(F)$ existe ssi $\mathcal{H}om(F,\mathcal{P}(e))$ existe, et ils sont isomorphes ; dans $\mathbf{Ens}$, $\mathcal{P}(e)$ a deux éléments.}
\entree{24265}{6:44:25}{Un sous-faisceau de $\mathcal{P}(F)$ est une famille d'ensembles de sous-faisceaux des $F|_U$, stable par restriction et de nature locale sur $U$.}
\entree{24395}{6:46:35}{Une topologie sur un faisceau $F$ se définit comme un sous-faisceau de $\mathcal{P}(F)$ vérifiant les axiomes usuels, lus dans $\mathrm{Sh}(X)$.}
\entree{24459}{6:47:39}{Pour $f:X\to Y$, il existe des flèches canoniques $f^*\mathcal{H}om(F,G)\to\mathcal{H}om(f^*F,f^*G)$ et $f^*\mathcal{P}(F)\to\mathcal{P}(f^*F)$.}
\entree{24649}{6:50:49}{Ces flèches ne sont pas en général des isomorphismes, par exemple pour un foncteur fibre ; elles le sont si $f$ est étale.}
\entree{24714}{6:51:54}{$f^*$ ne commute pas en général aux produits infinis ; il existe une flèche canonique $(\prod_i F_i)_x\to\prod_i(F_i)_x$.}
\entree{24851}{6:54:11}{Son injectivité signifie : si $f_i,g_i$ ont même germe en $x$ pour tout $i$, ils coïncident sur un même voisinage de $x$.}
\entree{24919}{6:55:19}{Contre-exemple : $U_i$ voisinages ouverts de $x$ d'intersection non voisinage, $F$ faisceau des applications vers $\{a,b\}$, $f_i=a$, $g_i=a$ sur $U_i$ et $b$ ailleurs.}
\entree{25045}{6:57:25}{Le foncteur fibre en $x$ commute aux produits ssi $x$ a un plus petit voisinage ouvert ; si les points sont fermés, ssi $x$ est isolé.}
\entree{25113}{6:58:33}{Dans $\mathrm{Sh}(X)$, un produit infini d'épimorphismes n'est pas toujours un épimorphisme ; on se ramène au cas des $F_i\to e$.}
\entree{25307}{7:01:47}{Exemple : $X$ localement connexe, $x$ fermé, $U_i$ voisinages ouverts de $x$ d'intersection non voisinage, $F_i=U_i\sqcup(X\setminus\{x\})$ ; $\prod_i F_i\to e$ n'est pas épi.}
\entree{25503}{7:05:03}{Si $X$ est compact totalement discontinu, $F\to G$ est épi ssi $F(U)\to G(U)$ est surjectif pour tout ouvert fermé $U$ ; tout produit d'épis est épi.}
\entree{25977}{7:12:57}{$X$ est connexe ssi l'objet final de $\mathrm{Sh}(X)$ n'est pas somme de deux sous-objets non initiaux.}
\face{Bande 2, face A}
\entree{26075}{7:14:35}{$X$ est non vide ssi $\varnothing\to e$ n'est pas un isomorphisme dans $\mathrm{Sh}(X)$ ; $X$ est vide ssi $\mathrm{Sh}(X)$ est équivalente à la catégorie ponctuelle.}
\entree{26202}{7:16:42}{Les objets groupes abéliens de $\mathrm{Sh}(X)$ forment une catégorie à assez d'injectifs ; $H^i(X,F)=R^i\Gamma(F)$, avec $\Gamma(F)=\mathrm{Hom}(e,F)$.}
\entree{26328}{7:18:48}{Pour un groupe abélien $M$, on pose $H^i(X,M)=H^i(X,M_X)$ ; ces groupes ne dépendent que de la catégorie $\mathrm{Sh}(X)$.}
\entree{26392}{7:19:52}{$X$ est connexe ssi $\mathbb{Z}\to H^0(X,\mathbb{Z})$ est surjectif ; connexe et non vide ssi il est bijectif.}
\entree{26457}{7:20:57}{Les sections de $I_X$ sur $U$ sont les applications localement constantes $U\to I$.}
\entree{26587}{7:23:07}{Le treillis des ouverts de $X$ est celui des sous-objets de l'objet final de $\mathrm{Sh}(X)$ ; les constructions de Čech (Artin--Mazur) s'en déduisent.}
\entree{26650}{7:24:10}{Si $X$ est connexe, localement connexe et localement simplement connexe, les revêtements de $X$ correspondent aux $\pi_1(X,x)$-ensembles.}
\entree{26716}{7:25:16}{Les revêtements de $X$ sont les objets localement constants $X'$ de $\mathrm{Sh}(X)$ : $X'\times U_i\simeq\coprod_{S_i}U_i$ pour un recouvrement $e=\sup_i U_i$.}
\entree{26781}{7:26:21}{Pour $X$ connexe et localement connexe mais non localement simplement connexe, les revêtements sont classifiés par un pro-groupe.}
\entree{26844}{7:27:24}{Pour un faisceau de groupes $G$ non nécessairement commutatif, $H^1(X,G)$ est l'ensemble des classes d'isomorphisme de $G$-torseurs.}
\entree{26911}{7:28:31}{Définition : un $G$-torseur est un objet $P$ muni d'une action de $G$ telle que $G\times P\to P\times P$ soit un isomorphisme et $P\to e$ un épimorphisme.}
\entree{26975}{7:29:35}{$G\mapsto G_X=\coprod_G e$ est exact à gauche ; les $G_X$-torseurs sont les revêtements galoisiens de groupe $G$.}
\entree{27040}{7:30:40}{Pour $X$ connexe, localement connexe et localement simplement connexe, $H^1(X,G)$ s'identifie à $\mathrm{Hom}(\pi_1(X,x),G)$ modulo conjugaison ; $X$ est simplement connexe ssi tous ces $H^1$ sont triviaux.}
\entree{27104}{7:31:44}{Définition : une catégorie fibrée sur $\mathcal{O}(X)$ associe à $U$ une catégorie $F_U$, à $V\subset U$ un foncteur de restriction, transitif à isomorphisme cohérent près.}
\entree{27167}{7:32:47}{Exemples (qui sont des champs) : $U\mapsto\mathrm{Sh}(U)$, espaces fibrés sur $U$, faisceaux de groupes sur $U$, $G|_U$-torseurs.}
\entree{27291}{7:34:51}{Définition : une catégorie fibrée est un champ si les préfaisceaux $\mathcal{H}om(F,G)$ sont des faisceaux et si toute donnée de recollement (cocycle sur les $U_{ijk}$) est effective.}
\entree{27549}{7:39:09}{Pour $G$ non commutatif, $H^2(X,G)$ se définit comme un ensemble de classes de champs (gerbes) attachés à $G$ (Giraud) ; il ne dépend que de $\mathrm{Sh}(X)$.}
\entree{27873}{7:44:33}{Définition : un champ sur $\mathcal{E}$ est une catégorie fibrée sur $\mathcal{E}$ avec descente le long des familles épimorphiques, les intersections devenant des produits fibrés.}
\entree{27938}{7:45:38}{Restriction et prolongement établissent une équivalence entre champs sur $\mathcal{O}(X)$ et champs sur $\mathrm{Sh}(X)$.}
\entree{28139}{7:48:59}{$\mathrm{Sh}(\coprod_i X_i)\simeq\prod_i\mathrm{Sh}(X_i)$ : la somme d'espaces correspond au produit des catégories.}
\entree{28203}{7:50:03}{Si $Z$ est fermé dans $X$ et dans $Y$, $\mathrm{Sh}(X\cup_Z Y)$ est le 2-produit fibré de $\mathrm{Sh}(X)$ et $\mathrm{Sh}(Y)$ au-dessus de $\mathrm{Sh}(Z)$.}
\entree{28331}{7:52:11}{Pour une relation d'équivalence $R$ sur $X$, les faisceaux sur $X$ munis d'une donnée de descente relative à $R$ forment un topos quotient, même si $X/R$ est pathologique.}
\entree{28459}{7:54:19}{Définition : un espace $Z$ est irréductible s'il est non vide et n'est pas réunion de deux fermés distincts de $Z$.}
\entree{28524}{7:55:24}{Un espace séparé est irréductible ssi c'est un point ; une variété est irréductible pour la topologie de Zariski ssi elle l'est au sens algébrique.}
\entree{28524}{7:55:24}{L'adhérence d'une partie irréductible est irréductible.}
\entree{28590}{7:56:30}{On note $\tilde X$ l'ensemble des fermés irréductibles de $X$ et $i:X\to\tilde X$, $x\mapsto\overline{\{x\}}$ ; $i$ est bijective si $X$ est séparé.}
\entree{28654}{7:57:34}{Les parties $\tilde U=\{Y\in\tilde X : Y\cap U\neq\varnothing\}$, $U$ ouvert, forment une topologie : $U\mapsto\tilde U$ respecte réunions et intersections finies.}
\entree{28850}{8:00:50}{On a $i^{-1}(\tilde U)=U$ ; $i$ est continue et induit un isomorphisme de treillis $\mathcal{O}(\tilde X)\simeq\mathcal{O}(X)$.}
\entree{28982}{8:03:02}{Par suite $i_*:\mathrm{Sh}(X)\to\mathrm{Sh}(\tilde X)$ est un isomorphisme de catégories, d'inverse $i^*$.}
\face{Bande 2, face B}
\entree{29036}{8:03:56}{$i:X\to\tilde X$ est injective ssi $X$ est $T_0$ ; si $X$ est séparé, $i$ est un homéomorphisme.}
\entree{29101}{8:05:01}{Définition : un point générique d'un fermé $Z$ est un point d'adhérence $Z$ ; $i$ est surjective ssi tout fermé irréductible a un point générique.}
\entree{29166}{8:06:06}{Exemple : une variété sur $K$ réduite à ses points $K$-rationnels a des fermés irréductibles sans point générique.}
\entree{29234}{8:07:14}{Définition : $X$ est sobre si tout fermé irréductible a exactement un point générique, c'est-à-dire si $i:X\to\tilde X$ est bijective, donc un homéomorphisme.}
\entree{29300}{8:08:20}{Pour $Y$ fermé irréductible, la fibre $F_Y=\varinjlim_{U\cap Y\neq\varnothing}F(U)$ (limite filtrante) est exacte à gauche et commute aux limites inductives.}
\entree{29427}{8:10:27}{$F_{\overline{\{x\}}}=F_x$ ; en général $F_Y$ est la fibre au point $Y$ de $F$ vu comme faisceau sur $\tilde X$.}
\entree{29490}{8:11:30}{$Y\mapsto(F\mapsto F_Y)$ est une bijection de $\tilde X$ sur les classes de foncteurs fibres $\mathrm{Sh}(X)\to\mathbf{Ens}$ (exacts à gauche, commutant aux $\varinjlim$).}
\entree{29559}{8:12:39}{Un foncteur fibre de $\mathrm{Sh}(X)$ n'a pas d'automorphisme non trivial.}
\entree{29622}{8:13:42}{$\tilde X$ est sobre, et $X\mapsto\tilde X$ est adjoint à gauche de l'inclusion des espaces sobres : $\mathrm{Hom}(X,Y)\simeq\mathrm{Hom}(\tilde X,Y)$ pour $Y$ sobre.}
\entree{29816}{8:16:56}{Sur l'ensemble $\Phi(\mathcal{E})$ des classes de foncteurs fibres, les $\tilde U=\{\varepsilon : \varepsilon(U)\neq\varnothing\}$, $U\subset e$, forment une topologie.}
\entree{30013}{8:20:13}{$X\to\Phi(\mathrm{Sh}(X))$ est un homéomorphisme ssi $X$ est sobre : un espace sobre se reconstruit à partir de $\mathrm{Sh}(X)$.}
\entree{30076}{8:21:16}{Si $X$ est $T_1$ (par exemple les points fermés d'un schéma de Jacobson), $X$ s'identifie au sous-espace des points fermés de $\tilde X$.}
\section{25 mai 1973}\label{s-1973-05-25}
\video{URU5jgTvjhU}
\noindent\textit{Vidéo 1/4 (\href{https://www.youtube.com/watch?v=URU5jgTvjhU}{\texttt{URU5jgTvjhU}}) : bande 1, face A (8:23:34--9:28:43); bande 1, face B (9:28:43--10:33:42); bande 2, face A (10:33:42--11:06:44).}\par\medskip
\face{Bande 1, face A}
\entree{30280}{8:24:40}{Pour $f:X\to Y$ continue, le point $f(x)$ correspond au foncteur fibre $\varepsilon_x\circ f^*$, car $(f^*F)_x\simeq F_{f(x)}$.}
\entree{30471}{8:27:51}{Si $X$ et $Y$ sont sobres, tout foncteur $\mathrm{Sh}(Y)\to\mathrm{Sh}(X)$ exact à gauche et commutant aux $\varinjlim$ est isomorphe à $f^*$ pour un unique $f$ continu.}
\entree{30599}{8:29:59}{Deux foncteurs sur $\mathrm{Sh}(Y)$ commutant aux limites inductives et isomorphes sur la sous-catégorie génératrice $\mathcal{O}(Y)$ sont isomorphes.}
\entree{30793}{8:33:13}{Pour $f,g:X\to Y$, un morphisme de $f$ vers $g$ est un morphisme entre $f^*$ et $g^*$ ; si $Y$ est séparé, il n'en existe que si $f=g$, et c'est l'identité.}
\entree{31066}{8:37:46}{Définition (première forme) : un topos est une catégorie ayant les propriétés d'exactitude de $\mathbf{Ens}$ ($\varinjlim$ quelconques, $\varprojlim$ finies) et un petit ensemble générateur.}
\entree{31132}{8:38:52}{Définition : $C\subset\mathrm{Ob}\,\mathcal{E}$ est générateur si $\mathcal{E}\to\widehat{C}$, $X\mapsto\mathrm{Hom}(-,X)|_C$, est fidèle.}
\entree{31258}{8:40:58}{Dans un topos, il revient au même que ce foncteur soit fidèle, pleinement fidèle, conservatif, ou conservatif sur les monomorphismes.}
\entree{31446}{8:44:06}{Définition (Giraud) : un topos est une catégorie avec (a) $\varprojlim$ finies, (b) sommes disjointes universelles, (c) relations d'équivalence effectives universelles, (d) petit générateur.}
\entree{31573}{8:46:13}{Définition : une somme $X=\coprod X_i$ est universelle si pour tout $X'\to X$, la flèche $\coprod_i(X_i\times_X X')\to X'$ est un isomorphisme.}
\entree{31704}{8:48:24}{Définition : elle est disjointe si les $X_i\to X$ sont des monomorphismes et si $X_i\times_X X_j$ est initial pour $i\neq j$.}
\entree{31896}{8:51:36}{Définition : une relation d'équivalence $R\subset X\times X$ (chaque $R(Z)$ en est une sur $X(Z)$) est effective si $R\simeq X\times_{X/R}X$.}
\entree{31962}{8:52:42}{Elle est universelle si $X\to X/R$ est un épimorphisme universel et si le passage au quotient commute aux changements de base $Y'\to X/R$.}
\entree{32152}{8:55:52}{Dans un topos, les sous-objets et les objets quotients d'un objet forment des ensembles petits.}
\entree{32217}{8:56:57}{$\mathcal{E}_{ab}$ est abélienne, à limites inductives filtrantes exactes et à générateur ; elle a donc assez d'injectifs.}
\entree{32283}{8:58:03}{Un objet injectif de $\mathcal{E}_{ab}$ est injectif dans $\mathcal{E}$ ; en particulier $\mathcal{E}$ a assez d'injectifs.}
\entree{32541}{9:02:21}{Dans un topos, $\Gamma(X)=\mathrm{Hom}(e,X)$ ; ses foncteurs dérivés sur $\mathcal{E}_{ab}$ sont les $H^i(\mathcal{E},-)$.}
\entree{32668}{9:04:28}{L'objet constant $I_{\mathcal{E}}=\coprod_I e$ : $I\mapsto I_{\mathcal{E}}$ commute aux $\varinjlim$ et aux $\varprojlim$ finies ; on pose $H^i(\mathcal{E},M)=H^i(\mathcal{E},M_{\mathcal{E}})$.}
\entree{32805}{9:06:45}{Un topos a des invariants d'homotopie (pro-groupes, méthode de Čech d'Artin--Mazur), qui redonnent les $\pi_i$ d'un espace localement contractile.}
\entree{32997}{9:09:57}{$\mathrm{Sh}(\varnothing)$ est équivalente à la catégorie ponctuelle : c'est le \q{topos vide}.}
\entree{33121}{9:12:01}{Un topos $\mathcal{E}$ est vide ssi la flèche $\varnothing_{\mathcal{E}}\to e_{\mathcal{E}}$ est un isomorphisme.}
\entree{33248}{9:14:08}{$\mathrm{Sh}(\mathrm{pt})\simeq\mathbf{Ens}$ par $F\mapsto F(\mathrm{pt})$ : c'est le topos ponctuel, ou \q{topos final}.}
\entree{33443}{9:17:23}{Un topos $\mathcal{E}$ est ponctuel ssi le foncteur $\mathbf{Ens}\to\mathcal{E}$, $I\mapsto I_{\mathcal{E}}$, est une équivalence.}
\entree{33505}{9:18:25}{Pour un groupe discret $G$ opérant sur $X$, la catégorie $\mathrm{Sh}(X,G)$ des espaces étalés sur $X$ munis d'une action compatible de $G$ est un topos.}
\entree{33698}{9:21:38}{$\mathrm{Sh}(X,\{1\})=\mathrm{Sh}(X)$ ; $\mathrm{Sh}(\mathrm{pt},G)=B_G$, catégorie des $G$-ensembles, est le \q{topos classifiant} de $G$.}
\entree{33822}{9:23:42}{$(B_G)_{ab}$ est la catégorie des $G$-modules, $\Gamma(M)=M^G$ et $H^i(B_G,M)=H^i(G,M)$, cohomologie de l'espace classifiant de $G$.}
\entree{33954}{9:25:54}{La cohomologie de $\mathrm{Sh}(X,G)$ est la cohomologie équivariante $H^i(X;G,M)$ (Tôhoku).}
\face{Bande 1, face B}
\entree{34130}{9:28:50}{Si $G\neq1$, le topos $B_G$ n'est équivalent à aucun $\mathrm{Sh}(X)$.}
\entree{34201}{9:30:01}{Le foncteur d'oubli $B_G\to\mathbf{Ens}$ est un foncteur fibre, de groupe d'automorphismes $G$.}
\entree{34328}{9:32:08}{Si $X$ est sobre, tout foncteur fibre de $\mathrm{Sh}(X)$ est isomorphe à $F\mapsto F_x$ pour un unique $x\in X$.}
\entree{34393}{9:33:13}{La spécialisation ($y\in\overline{\{x\}}$) est un préordre, qui est un ordre si $X$ est sobre ; elle induit $F_y\to F_x$, fonctoriel en $F$.}
\entree{34662}{9:37:42}{Il existe un morphisme $\varepsilon_y\to\varepsilon_x$ ssi $y$ est une spécialisation de $x$, et il est unique : les points de $\mathrm{Sh}(X)$ forment un ensemble ordonné.}
\entree{34728}{9:38:48}{$B_G$ a un seul point à isomorphisme près sans être le topos final : ses points ont des automorphismes non triviaux.}
\entree{34858}{9:40:58}{Dans $\mathrm{Sh}(X,G)$, $\varepsilon^G_x:F\mapsto F_x$ est un point, de groupe d'automorphismes le stabilisateur $G_x$.}
\entree{35193}{9:46:33}{Théorème (Giraud) : $\mathcal{E}$ est un topos ssi il existe $C$ petite et $i:\mathcal{E}\to\widehat{C}$ pleinement fidèle ayant un adjoint à gauche exact à gauche.}
\entree{35463}{9:51:03}{Pour toute sous-catégorie pleine génératrice $C$ d'un topos, $\mathcal{E}\to\widehat{C}$ est pleinement fidèle et a un adjoint à gauche exact à gauche.}
\entree{35591}{9:53:11}{Les topos au sens des axiomes de Giraud sont exactement les catégories de faisceaux sur un site.}
\entree{35724}{9:55:24}{Dans un topos, $\mathcal{H}om(X,Y)$ existe : il représente $Z\mapsto\mathrm{Hom}(X\times Z,Y)$.}
\entree{35789}{9:56:29}{Dans un topos, le foncteur $Z\mapsto\mathrm{Sub}(Z)$ est représentable par un objet $P_{\mathcal{E}}$.}
\entree{35860}{9:57:40}{Lemme : dans un topos, $P:\mathcal{E}^{\circ}\to\mathbf{Ens}$ est représentable ssi il transforme limites inductives en limites projectives.}
\entree{36055}{10:00:55}{Tout topos est un topos élémentaire (Lawvere--Tierney : $\varprojlim$ finies, $\mathcal{H}om$, $P_{\mathcal{E}}$) ; un topos élémentaire à sommes petites et petit générateur est un topos.}
\entree{36307}{10:05:07}{Dans $\mathrm{Sh}(X,G)$, limites, $\mathcal{H}om$ et $P$ se calculent dans $\mathrm{Sh}(X)$, avec l'action induite de $G$.}
\entree{36436}{10:07:16}{Pour $C$ petite, $\widehat{C}$ est un topos (limites argument par argument, représentables générateurs), avec toutes les propriétés d'exactitude de $\mathbf{Ens}$.}
\entree{36630}{10:10:30}{$B_G=\widehat{C}$ pour $C$ la catégorie à un objet d'endomorphismes $G$.}
\entree{36693}{10:11:33}{Pour $X\in C$, $F\mapsto F(X)$ est un point de $\widehat{C}$ ; tout point est limite inductive filtrante de ceux-là, et $\widehat{C}$ a assez de points.}
\entree{36756}{10:12:36}{$\mathrm{Sh}(X,G)$ a assez de points : les $\varepsilon^G_x$ forment une famille conservative.}
\entree{36882}{10:14:42}{Pour un objet groupe $G$ d'un topos $\mathcal{E}$, les objets de $\mathcal{E}$ munis d'une action de $G$ forment un topos $B_G$ ; pour $\mathcal{E}=\mathbf{Ens}$, les $G$-ensembles.}
\entree{37075}{10:17:55}{$H^i(G,M)=H^i(B_G,M)$ pour $M$ abélien muni d'une action de $G$ ; $\Gamma(X)$ est l'ensemble des sections de $X$ invariantes par $G$.}
\entree{37271}{10:21:11}{$B_G$ classifie les $G$-torseurs : pour $f:\mathcal{T}\to\mathcal{E}$, les morphismes $\mathcal{T}\to B_G$ au-dessus de $\mathcal{E}$ correspondent aux $f^*G$-torseurs.}
\entree{37400}{10:23:20}{Pour un objet $X$ d'un topos $\mathcal{E}$, $\mathcal{E}/X$ est un topos, le topos induit ; $\mathrm{Sh}(Y)/X\simeq\mathrm{Sh}(X)$ pour $X$ étale sur $Y$.}
\entree{37930}{10:32:10}{Entre topos, tout foncteur $f^*$ commutant aux limites inductives a un adjoint à droite $f_*$ (lemme de représentabilité).}
\face{Bande 2, face A}
\entree{38154}{10:35:54}{Définition : un morphisme de topos $f:\mathcal{E}\to\mathcal{F}$ est un couple de foncteurs adjoints $f^*\dashv f_*$, avec $f^*:\mathcal{F}\to\mathcal{E}$ exact à gauche.}
\entree{38285}{10:38:05}{Les morphismes $\mathcal{E}\to\mathcal{F}$ forment une catégorie $\mathbf{Homtop}(\mathcal{E},\mathcal{F})$, de flèches $f\to g$ les morphismes $f_*\to g_*$ (ou $g^*\to f^*$).}
\entree{38728}{10:45:28}{Composition : $(gf)_*=g_*f_*$ et $(gf)^*=f^*g^*$ ; elle est associative.}
\entree{38859}{10:47:39}{La composition est un foncteur $\mathbf{Homtop}(\mathcal{E},\mathcal{F})\times\mathbf{Homtop}(\mathcal{F},\mathcal{G})\to\mathbf{Homtop}(\mathcal{E},\mathcal{G})$ : les topos forment une 2-catégorie.}
\entree{39237}{10:53:57}{Un morphisme $f$ est une équivalence de topos ssi $f^*$, ou $f_*$, est une équivalence de catégories.}
\entree{39560}{10:59:20}{Pour $Y$ sobre, la catégorie des points $\mathbf{Homtop}(\mathrm{pt},\mathrm{Sh}(Y))$ est équivalente à $Y$ ordonné par spécialisation.}
\entree{39689}{11:01:29}{Pour $X,Y$ sobres, $\mathbf{Homtop}(\mathrm{Sh}(X),\mathrm{Sh}(Y))$ est équivalente aux $f$ continues ordonnées par $f\le g$ ssi $g(x)\in\overline{\{f(x)\}}$ pour tout $x$.}
\entree{39885}{11:04:45}{Si $Y$ est séparé, cette catégorie est discrète.}
\section{30 mai 1973}\label{s-1973-05-30}
\video{PqCQM3Pcl6Q}
\noindent\textit{Vidéo 2/4 (\href{https://www.youtube.com/watch?v=PqCQM3Pcl6Q}{\texttt{PqCQM3Pcl6Q}}) : bande 1, face A (0:00:00--0:49:07); bande 1, face B (0:49:07--1:38:30); bande 2, face A (1:38:30--1:55:37).}\par\medskip
\noindent \emph{À venir} : la transcription de cette séance n'est pas encore faite.\par
\section{4 juin 1973}\label{s-1973-06-04}
\video{PqCQM3Pcl6Q}
\noindent\textit{Vidéo 2/4 (\href{https://www.youtube.com/watch?v=PqCQM3Pcl6Q}{\texttt{PqCQM3Pcl6Q}}) : bande 1, face A (1:55:37--2:44:38); bande 1, face B (2:44:38--3:33:35); bande 2 (3:33:35--4:22:54).}\par\medskip
\noindent \emph{À venir} : la transcription de cette séance n'est pas encore faite.\par
\section{6 juin 1973}\label{s-1973-06-06}
\video{PqCQM3Pcl6Q}
\noindent\textit{Vidéo 2/4 (\href{https://www.youtube.com/watch?v=PqCQM3Pcl6Q}{\texttt{PqCQM3Pcl6Q}}) : bande 1 (4:22:54--5:12:09).}\par\medskip
\noindent \emph{À venir} : la transcription de cette séance n'est pas encore faite.\par
\section{20 juin 1973}\label{s-1973-06-20}
\video{PqCQM3Pcl6Q}
\noindent\textit{Vidéo 2/4 (\href{https://www.youtube.com/watch?v=PqCQM3Pcl6Q}{\texttt{PqCQM3Pcl6Q}}) : bande 1, face A (5:12:09--6:01:04); bande 1, face B (6:01:04--6:50:09); bande 2, face A (6:50:09--7:38:43); bande 2, face B (7:38:43--8:27:21).}\par\medskip
\noindent \emph{À venir} : la transcription de cette séance n'est pas encore faite.\par
\section{27 juin 1973}\label{s-1973-06-27}
\video{PqCQM3Pcl6Q}
\noindent\textit{Vidéo 2/4 (\href{https://www.youtube.com/watch?v=PqCQM3Pcl6Q}{\texttt{PqCQM3Pcl6Q}}) : bande 1, face A (8:27:21--9:02:23); bande 2, face A (9:02:23--9:52:24); bande 2, face B (9:52:24--10:14:40).}\par\medskip
\noindent \emph{À venir} : la transcription de cette séance n'est pas encore faite.\par
\section{2 juillet 1973}\label{s-1973-07-02}
\video{4gynb9uTEE4}
\noindent\textit{Vidéo 3/4 (\href{https://www.youtube.com/watch?v=4gynb9uTEE4}{\texttt{4gynb9uTEE4}}) : bande 1 (0:00:00--0:48:52); bande 2 (0:48:52--1:54:02); bande 3, face A (1:54:02--2:43:03); bande 3, face B (2:43:03--3:31:57).}\par\medskip
\face{Bande 1}
\entree{258}{0:04:18}{Pour $\Delta$ = catégories à limites projectives quelconques, la localisation $C\to C[\Sigma^{-1}]$ prise dans $\Delta$ est bijective sur les objets : ajouter des axiomes ajoute des flèches, non des constructions.}
\entree{321}{0:05:21}{Une catégorie $I$ est strictement cogénératrice dans son enveloppe $\widetilde{I}$ par limites projectives quelconques.}
\entree{386}{0:06:26}{Définition : $J\subset E$ est \q{strictement génératrice} si tout objet de $E$ est limite inductive d'un diagramme d'objets de $J$ ; exemple : $J$ dans $\widehat{J}$.}
\entree{451}{0:07:31}{Si $I$ est strictement cogénératrice dans chaque $R_\alpha$ d'un système filtrant de $\Delta$, elle l'est dans sa 2-limite inductive prise dans $\Delta$.}
\entree{648}{0:10:48}{La stricte cogénération est transitive ; elle se conserve par adjonction de flèches au sens de $\Delta$ et par localisation dans $\Delta$.}
\entree{712}{0:11:52}{Corollaire : à chaque étape de la construction d'une théorie algébrique de base $I$, l'image de $I$ reste strictement cogénératrice.}
\entree{776}{0:12:56}{Théorème : pour $\Delta$ = catégories à limites projectives quelconques, les théories algébriques sont celles représentées par un $R\in\Delta$ ayant une petite sous-catégorie strictement cogénératrice.}
\entree{848}{0:14:08}{De plus, $I\to R$ est basique si et seulement si l'image de $I$ est strictement cogénératrice ; dualement, $R^{\mathrm{op}}$ a des limites inductives quelconques et un petit ensemble strictement générateur.}
\entree{1104}{0:18:24}{Si $I'\subset R$ est une petite sous-catégorie pleine strictement cogénératrice, le $\Delta$-foncteur $\widetilde{I'}\to R$ est essentiellement surjectif et plein ; $R$ est une catégorie de fractions de $\widetilde{I'}$.}
\entree{1357}{0:22:37}{Corollaire : une telle théorie s'obtient en trois étapes (structure vide sur $I$, adjonction de flèches, axiomes), sans construction transfinie.}
\entree{1845}{0:30:45}{Si $T$ est représentée par $R$ ayant une petite sous-catégorie strictement cogénératrice, $T(\mathbf{Ens})=\mathrm{Hom}_\Delta(R,\mathbf{Ens})\simeq R^{\mathrm{op}}$ : $T$ est déterminée par ses réalisations ensemblistes.}
\entree{2233}{0:37:13}{Exemple : pour la théorie des groupes, $R\simeq\mathbf{Grp}^{\mathrm{op}}$ ; l'objet universel est $\mathbb{Z}$, cogroupe de $\mathbf{Grp}$ par $\mathbb{Z}\to\mathbb{Z}*\mathbb{Z}$, $x\mapsto x_1x_2$.}
\entree{2681}{0:44:41}{Pour $C$ petite et une classe $D'$ de diagrammes, l'enveloppe de $C$ par limites de type $D'$ est la plus petite sous-catégorie pleine de l'enveloppe $\widetilde{C}$ contenant $C$ et stable par ces limites.}
\entree{2746}{0:45:46}{Exemple : sur les catégories à produits finis, la théorie des groupes est représentée par une petite catégorie, la sous-catégorie pleine de $\mathbf{Grp}^{\mathrm{op}}$ des groupes libres de type fini.}
\entree{2813}{0:46:53}{Une théorie sur les catégories à limites projectives de type $D'$ en définit par restriction une sur les catégories à limites projectives quelconques ; il n'y a pas d'opération inverse en général.}
\face{Bande 2}
\entree{3075}{0:51:15}{Si $F:E\to E'$ est $\pi$-accessible et envoie $E_\pi$ dans $E'_\pi$ (objets $\pi$-accessibles), $F$ est le prolongement $\mathrm{Ind}_\pi$ de sa restriction $E_\pi\to E'_\pi$.}
\entree{3135}{0:52:15}{Pour $F$ accessible et $\pi$ assez grand, $F(E_\pi)\subset E'_{\pi^{\pi_0}}$ ($\pi_0$ fixé, SGA 4) ; comme $(2^\lambda)^{\pi_0}=2^\lambda$ pour $\lambda\ge\pi_0$ infini, $F(E_\pi)\subset E'_\pi$ si $\pi=2^\lambda$.}
\entree{3263}{0:54:23}{Corollaire : un foncteur accessible entre catégories accessibles respecte la filtration canonique par les objets $\pi$-accessibles pour des $\pi$ arbitrairement grands.}
\entree{3398}{0:56:38}{Pour $i:C\hookrightarrow E$ petite pleine et $E'$ à limites inductives, la restriction $i^*:\mathrm{Hom}(E,E')\to\mathrm{Hom}(C,E')$ a un adjoint à gauche $i_!$ pleinement fidèle.}
\entree{3460}{0:57:40}{Pour $E$ accessible, $F:E\to E'$ est accessible si et seulement s'il est dans l'image essentielle de $i_!$ pour une petite $C\subset E$.}
\entree{3586}{0:59:46}{Pour $C$ petite à idempotents scindés, $\mathrm{Ind}_\pi(C)$ a des limites inductives quelconques si et seulement si $C$ a les limites inductives de cardinal au plus $\pi$.}
\entree{3791}{1:03:11}{Pour une telle $C$, $\mathrm{Ind}_\pi(C)$ est l'enveloppe $\Delta_\pi$--$\Delta$ de $C$ : $\mathrm{Hom}_\Delta(\mathrm{Ind}_\pi C,F)\simeq\mathrm{Hom}_{\Delta_\pi}(C,F)$ pour toute $F$ à limites inductives quelconques.}
\entree{3986}{1:06:26}{Les catégories $R$ représentant les théories algébriques sont exactement les catégories \q{co-accessibles} : $S=T(\mathbf{Ens})\simeq R^{\mathrm{op}}$ est accessible (et cocomplète).}
\entree{4186}{1:09:46}{Sur une catégorie accessible cocomplète $E$, un foncteur $E^{\mathrm{op}}\to\mathbf{Ens}$ est représentable si et seulement s'il transforme limites inductives en limites projectives.}
\entree{4249}{1:10:49}{Un foncteur $F:E\to\mathbf{Ens}$ ($E$ accessible cocomplète) est représentable si et seulement s'il commute aux limites projectives quelconques et est accessible.}
\entree{4326}{1:12:06}{SGA 4 : un foncteur $F:E\to\mathbf{Ens}$ est pro-représentable par un pro-objet à petit indexage si et seulement s'il est exact à gauche et accessible.}
\entree{4391}{1:13:11}{Si $E$ a des limites projectives quelconques, un foncteur pro-représentable $E\to\mathbf{Ens}$ est représentable si et seulement s'il commute à ces limites.}
\entree{4454}{1:14:14}{Entre catégories accessibles cocomplètes, $F$ a un adjoint à droite si et seulement s'il commute aux limites inductives, un adjoint à gauche si et seulement s'il commute aux limites projectives et est accessible.}
\entree{4531}{1:15:31}{Exemple (SGA 4) : sur $\mathbf{Grp}$, un foncteur à petites valeurs, pro-représenté par un système non petit formé de groupes simples, commute aux limites projectives sans être accessible.}
\entree{4664}{1:17:44}{Un morphisme de théories algébriques $T\to T'$ équivaut à un foncteur $S'\to S$ commutant aux limites inductives ; son adjoint à droite $S\to S'$ est $T(\mathbf{Ens})\to T'(\mathbf{Ens})$.}
\entree{4853}{1:20:53}{Les foncteurs $S\to S'$ provenant de morphismes de théories algébriques sont exactement ceux qui commutent aux limites projectives et sont accessibles.}
\entree{4931}{1:22:11}{Si $F\to I$ est fibrée sur $I$ petite, à fibres et foncteurs de transition accessibles, la catégorie des sections cartésiennes (2-limite projective) et celle de toutes les sections sont accessibles.}
\entree{5061}{1:24:21}{Corollaire : pour $E$ accessible et $I$ petite, $\mathrm{Hom}(I,E)$ est accessible, ainsi que $\mathrm{Hom}(I[I^{-1}],E)$, foncteurs transformant toute flèche en isomorphisme.}
\entree{5320}{1:28:40}{La 2-limite inductive d'un petit système de catégories accessibles cocomplètes, prise dans les catégories cocomplètes et foncteurs cocontinus, est accessible.}
\entree{5463}{1:31:03}{Question : la 2-limite inductive dans $\mathbf{Cat}$ d'un petit système de catégories accessibles est-elle accessible ?}
\entree{5529}{1:32:09}{Exemple : la somme dans $\mathbf{Cat}$ de deux catégories non vides à limites inductives n'a pas d'objet initial, ni de somme d'un objet de chacune.}
\entree{6306}{1:45:06}{Pour un morphisme de théories algébriques, $T(\mathbf{Ens})\to T'(\mathbf{Ens})$ a un adjoint à gauche (structure libre) ; exemple : groupe enveloppant d'un monoïde.}
\entree{6434}{1:47:14}{Si un foncteur pleinement fidèle admet un adjoint à gauche, cet adjoint est une localisation (passage à une catégorie de fractions).}
\entree{6499}{1:48:19}{Les sous-théories pleines de $T_{R'}$ correspondent aux localisations $R'\to R'[\Sigma'^{-1}]$ dans $\Delta$, $\Sigma'$ stable par limites projectives, et aux sous-catégories pleines réflexives de $S'$ stables par limites projectives.}
\entree{6627}{1:50:27}{Pour une sous-théorie algébrique, $\Sigma'$ est engendré, par opérations de limites projectives, par un petit ensemble de flèches.}
\entree{6691}{1:51:31}{Question : une sous-catégorie pleine $S$ d'une catégorie accessible $S'$, stable par limites projectives et d'inclusion admettant un adjoint à gauche, est-elle accessible ?}
\face{Bande 3, face A}
\entree{7050}{1:57:30}{Clause de passage à la limite : la 2-limite projective, au-dessus de $T^I$, d'un système filtrant de théories algébriques de base $I$ est algébrique de base $I$.}
\entree{7114}{1:58:34}{Les théories représentables forment une 2-sous-catégorie pleine des $\Delta$-théories, 2-équivalente à $\Delta^{\mathrm{op}}$ par $R\mapsto T_R$ ; elle contient les théories algébriques.}
\entree{7240}{2:00:40}{$R$ est algébrique de base $I$ (discrète) s'il est le terme d'une suite transfinie de $\Delta$ partant de l'enveloppe de $I$, chaque pas inversant, égalisant ou ajoutant des flèches, avec 2-limites inductives aux ordinaux limites.}
\entree{7434}{2:03:54}{La fidélité de $\mathrm{Hom}_\Delta(R,C)\to C^I$ pour tout $C$ ne suffit pas pour que $I\to R$ soit basique (contre-exemple connu).}
\entree{7434}{2:03:54}{\q{Problème trois} (ouvert) : $I\to R$ est-il basique si et seulement si toute sous-catégorie strictement pleine de $R$ stable par limites de type $D$ et contenant l'image de $I$ est égale à $R$ ?}
\entree{7564}{2:06:04}{Cette condition équivaut à : la saturation transfinie de l'image de $I$ par limites de type $D'$ et $D''$ atteint $R$.}
\entree{7694}{2:08:14}{Si $D'$ et $D''$ sont essentiellement petits (et si les problèmes un et deux conservent la petitesse), les objets algébriques de $\Delta$ sont exactement les petits.}
\entree{7823}{2:10:23}{Énoncé conjectural : si $R$ est algébrique, on peut agrandir $I$ pour n'avoir plus de flèches à ajouter ; $T_R$ est alors une sous-théorie pleine de $C\mapsto\mathrm{Hom}(I',C)$.}
\entree{8276}{2:17:56}{Condition nécessaire de 2-représentabilité : $T$ commute aux 2-limites projectives, $T(\varprojlim C_i)\simeq\varprojlim T(C_i)$.}
\entree{8339}{2:18:59}{Sur une catégorie à limites projectives quelconques, un foncteur vers $\mathbf{Ens}$ qui commute à ces limites est pro-représentable par un système projectif éventuellement non petit.}
\entree{8410}{2:20:10}{Si $T$ est 2-représentable, $T$ transforme les $\Delta$-foncteurs fidèles (resp. pleinement fidèles) en foncteurs fidèles (resp. pleinement fidèles).}
\entree{8607}{2:23:27}{Propriété d'exactitude donnée par $R\in\Delta_D$ et $\Sigma$ : tout $\Delta_D$-foncteur $R\to C$ inverse $\Sigma$ ; équivaut à : $\mathrm{Hom}_{\Delta_D}(R',C)\to\mathrm{Hom}_{\Delta_D}(R,C)$ est une équivalence, $R'$ localisé.}
\entree{8738}{2:25:38}{\q{Propriétés d'exactitude généralisées} : pour un $\Delta_D$-foncteur $R\to R'$ donné, exiger que $\mathrm{Hom}(R',C)\to\mathrm{Hom}(R,C)$ soit une équivalence, pleinement fidèle ou fidèle.}
\entree{8802}{2:26:42}{Les sous-2-catégories de $\Delta_D$ ainsi définies sont stables par 2-limites projectives dans $\mathbf{Cat}$.}
\entree{8931}{2:28:51}{Exemple ($D=\varnothing$) : imposer l'inversion de la flèche de $\{0\to1\}$ définit les groupoïdes.}
\entree{8998}{2:29:58}{Égaliser les deux flèches de $\{0\rightrightarrows1\}$ définit les catégories préordonnées, équivalentes aux ensembles préordonnés, 2-équivalentes aux ensembles ordonnés.}
\entree{9126}{2:32:06}{Les deux conditions ensemble définissent les catégories essentiellement discrètes, 2-équivalentes à $\mathbf{Ens}$.}
\entree{9190}{2:33:10}{Pour $\Delta=\mathbf{Cat}$ : l'enveloppe de $I$ est $I$ ; le problème deux est résolu par la catégorie de fractions $I[\Sigma^{-1}]$ (Gabriel--Zisman).}
\entree{9256}{2:34:16}{Pour les groupoïdes : l'enveloppe de $I$ est son groupoïde fondamental $I[I^{-1}]$, de mêmes composantes connexes ; le problème deux est trivial.}
\entree{9381}{2:36:21}{Pour les préordres : l'enveloppe de $I$ est $\mathrm{Ob}\,I$ muni du préordre $x\ge y$ s'il existe une flèche $x\to y$.}
\entree{9444}{2:37:24}{Rendre inversibles des relations $x_i\le y_i$ d'un préordre : prendre le préordre engendré par le préordre donné et les $y_i\le x_i$.}
\entree{9510}{2:38:30}{Pour les catégories discrètes : l'enveloppe de $I$ est $\pi_0(I)$, ensemble des composantes connexes ; le problème deux est trivial.}
\entree{9714}{2:41:54}{Si $T=T_R$ et si $C$ a des limites projectives de type $\delta$ commutant aux limites inductives de type $D''$, $T(C)$ a ces limites, calculées argument par argument.}
\face{Bande 3, face B}
\entree{9915}{2:45:15}{Un morphisme de théories représentables $T\to T'$ donne $T(C)\to T'(C)$ commutant à ces limites ; de même $T(C)\to T(C')$ si $C\to C'$ les préserve.}
\entree{9978}{2:46:18}{Pour tout objet $r$ de $R$, le foncteur d'évaluation $r_C:T(C)\to C$ commute aux limites projectives de type $\delta$.}
\entree{10040}{2:47:20}{Exemple : la catégorie des torseurs sous un groupe $G\neq1$ n'a ni produits, ni sommes, ni objet initial ou final.}
\entree{10103}{2:48:23}{Si $D''=\varnothing$ et si $C$ a des limites projectives de type $\delta$ (resp. quelconques), $T(C)$ en a aussi.}
\entree{10240}{2:50:40}{Pour $\Delta'\subset\Delta$ (plus de limites exigées), la restriction envoie $\Delta$-théories sur $\Delta'$-théories et théories algébriques sur théories algébriques.}
\entree{10370}{2:52:50}{Si les problèmes un et deux ont une solution dans $\Delta'$, la restriction de $T_R$ à $\Delta'$ est représentée par l'\q{enveloppe $\Delta$--$\Delta'$} $R'$ de $R$.}
\entree{10564}{2:56:04}{$R'$ s'obtient de l'enveloppe de $R$ dans $\Delta'$ en rendant inversibles les flèches de comparaison des limites de type $D$ de $R$.}
\entree{10915}{3:01:55}{Le foncteur $\mathrm{Hom}(T_1,T_2)\to\mathrm{Hom}(T'_1,T'_2)$ de restriction n'est en général ni fidèle, ni plein, ni essentiellement surjectif.}
\entree{11448}{3:10:48}{Pour $J$ petite, $J\to\widehat{J}$ est l'enveloppe par limites inductives : pour $D$ cocomplète, la restriction des foncteurs cocontinus $\widehat{J}\to D$ à $J$ est une équivalence.}
\entree{11509}{3:11:49}{Tout préfaisceau $F$ sur $J$ est limite inductive sur $J/F$ des objets de $J$ au-dessus de $F$ ; le prolongement est $\widehat{\varphi}(F)=\varinjlim_{J/F}\varphi(X)$.}
\entree{11699}{3:14:59}{Problème un pour $\Delta$ = catégories à petites limites projectives : l'enveloppe de $I$ petite est $\mathrm{Hom}(I,\mathbf{Ens})^{\mathrm{op}}$, avec le plongement de Yoneda.}
\entree{11824}{3:17:04}{Dans $\Delta$, égaliser $f,g:X\rightrightarrows Y$ revient à rendre inversible le noyau $\mathrm{Ker}(f,g)\to X$.}
\entree{11886}{3:18:06}{Problème deux : si $\Sigma'$ est le saturé de $\Sigma$, la catégorie de fractions $C[\Sigma'^{-1}]$ prise dans $\mathbf{Cat}$ est dans $\Delta$, et $C\to C[\Sigma'^{-1}]$ est la solution.}
\entree{12093}{3:21:33}{Le saturé $\Sigma'$ (flèches inversées par tout $\Delta$-foncteur inversant $\Sigma$) contient les isomorphismes, est stable par composition et changement de base, et si deux de $u$, $v$, $vu$ y sont, le troisième aussi.}
\entree{12219}{3:23:39}{Si $u:X\to Y$ est dans $\Sigma'$, la diagonale $X\to X\times_Y X$ est dans $\Sigma'$.}
\entree{12284}{3:24:44}{$\Sigma'$ admet un calcul de fractions : $\mathrm{Hom}(X,Y)=\varinjlim\mathrm{Hom}_C(X',Y)$, sur le système projectif filtrant des $X'\to X$ dans $\Sigma'$.}
\entree{12473}{3:27:53}{Si $C$ a des limites projectives finies, $C[\Sigma'^{-1}]$ aussi, et le foncteur de localisation est exact à gauche.}
\entree{12539}{3:28:59}{$\Sigma'$ est stable par produits fibrés quelconques : si les $X_i\to Y$ sont dans $\Sigma'$, $\prod_Y X_i\to Y$ aussi.}
\entree{12667}{3:31:07}{$C[\Sigma'^{-1}]$ a des produits quelconques, images de ceux de $C$ ; ainsi $C\to C[\Sigma'^{-1}]$ est un $\Delta$-foncteur.}
\section{10 juillet 1973}\label{s-1973-07-10}
\video{4gynb9uTEE4}
\noindent\textit{Vidéo 3/4 (\href{https://www.youtube.com/watch?v=4gynb9uTEE4}{\texttt{4gynb9uTEE4}}) : bande 1, face A (3:31:57--4:36:29); bande 1, face B (4:36:29--5:41:33).}\par\medskip
\face{Bande 1, face A}
\entree{12928}{3:35:28}{Pour $E$ dans $\Delta$ et une classe saturée $\widetilde\Sigma$ non petite, la catégorie de fractions $E[\widetilde\Sigma^{-1}]$ peut avoir des $\mathrm{Hom}$ non petits, et n'est alors pas dans $\Delta$.}
\entree{13056}{3:37:36}{Pour une théorie algébrique $R$ d'objets de base $I$, $I$ n'est pas en général strictement cogénérateur dans $R$ : une structure de groupe n'est pas déterminée par les $x\mapsto x^n$.}
\entree{13185}{3:39:45}{La stricte génération n'est pas transitive : $\mathbb{Z}$ est strictement générateur dans les groupes libres, ceux-ci le sont dans $\mathbf{Gr}$, mais $\{\mathbb{Z}\}$ ne l'est pas dans $\mathbf{Gr}$ (de même dans $\mathbf{Ab}$).}
\entree{13314}{3:41:54}{Une théorie algébrique relative aux limites projectives ne s'obtient pas toujours en deux étapes (flèches, puis axiomes) ; il peut falloir une construction transfinie, indexée par un petit ordinal.}
\entree{13380}{3:43:00}{Si $R\in\Delta$ est algébrique, $R$ admet une petite famille strictement cogénératrice ; la réciproque est fausse.}
\entree{13448}{3:44:08}{Exemple : $\mathbf{Ens}^{\mathrm{op}}$ admet un générateur strict (l'ensemble à deux éléments) mais n'est la catégorie $T(\mathbf{Ens})$ d'aucune théorie algébrique $T$.}
\entree{13577}{3:46:17}{$\mathbf{Ens}^{\mathrm{op}}$ est équivalente à la catégorie des algèbres de Boole complètes atomiques, par $X\mapsto\mathcal{P}(X)$.}
\entree{13967}{3:52:47}{La complétude d'une algèbre de Boole porte sur des familles indexées par un ensemble $I$ arbitraire : famille non petite d'opérations et d'axiomes, exclue des théories algébriques du cours.}
\entree{14825}{4:07:05}{Toute algèbre de Boole complète relève, pour un cardinal $\pi$ convenable, de la théorie des algèbres de Boole $\pi$-complètes, qui n'a qu'un petit ensemble d'opérations.}
\entree{15016}{4:10:16}{Les objets de $\Delta$ représentant des théories algébriques sont les $\Delta$-enveloppes (adjonction des petites limites projectives) de petites catégories de $\Delta_\pi$ (limites projectives de cardinal ${\le}\,\pi$).}
\entree{15081}{4:11:21}{Les catégories $S=R^{\mathrm{op}}=T(\mathbf{Ens})$ des réalisations ensemblistes de ces théories sont des catégories accessibles.}
\entree{15284}{4:14:44}{Définition : un ensemble préordonné $I$ \q{domine $\pi$} s'il est filtrant et si toute famille de cardinal ${\le}\,\pi$ d'éléments de $I$ est majorée ; pour $\pi\ge 2$, la seconde condition entraîne la première.}
\entree{15352}{4:15:52}{Définition : $E$ vérifie $L_\pi$ si les limites inductives indexées par tout petit $I$ dominant $\pi$ existent dans $E$ ; $E$ vérifie $L$ si elle vérifie $L_\pi$ pour un petit cardinal $\pi$.}
\entree{15482}{4:18:02}{Définition : $F\colon E\to E'$ est $\pi$-accessible si $E$ vérifie $L_\pi$ et si $F$ commute aux limites inductives indexées par les $I$ dominant $\pi$ ; accessible s'il est $\pi$-accessible pour un petit $\pi$.}
\entree{15673}{4:21:13}{Définition : un objet $X$ de $E$ est $\pi$-accessible (resp. accessible) si le foncteur $\mathrm{Hom}(X,-)$ de $E$ dans $\mathbf{Ens}$ est $\pi$-accessible (resp. accessible).}
\entree{15809}{4:23:29}{Définition : $E$ est $\pi$-accessible si elle a une petite sous-catégorie pleine strictement génératrice $C$, formée d'objets $\pi$-accessibles (tout objet est limite inductive d'un diagramme de $C$).}
\entree{15877}{4:24:37}{Cette définition est plus forte que celle de SGA 4, I, 9, qui ne demandait qu'une petite sous-catégorie génératrice.}
\entree{16008}{4:26:48}{Si $E$, $F$ vérifient $L_\pi$, la sous-catégorie pleine $\mathrm{Hom}_\pi(E,F)$ des foncteurs $\pi$-accessibles est stable dans $\mathrm{Hom}(E,F)$ par limites inductives calculées argument par argument.}
\entree{16076}{4:27:56}{Elle est stable aussi par les limites projectives, calculées argument par argument, indexées par une petite catégorie $J$ ayant au plus $\pi$ flèches.}
\entree{16213}{4:30:13}{Dans $\mathbf{Ens}$, les limites inductives suivant un $I$ dominant $\pi$ commutent aux limites projectives suivant une catégorie ayant au plus $\pi$ flèches.}
\entree{16213}{4:30:13}{Les foncteurs accessibles forment une sous-catégorie pleine de $\mathrm{Hom}(E,F)$ stable par petites limites inductives et projectives calculées argument par argument.}
\entree{16276}{4:31:16}{Si $E$ vérifie $L_\pi$, la sous-catégorie pleine $E_\pi$ des objets $\pi$-accessibles est stable par limites inductives indexées par une catégorie ayant au plus $\pi$ flèches.}
\entree{16341}{4:32:21}{Si $E$ vérifie $L$, la sous-catégorie pleine des objets accessibles de $E$ est stable par petites limites inductives.}
\entree{16406}{4:33:26}{Si $E$ est accessible, tout objet de $E$ est accessible.}
\face{Bande 1, face B}
\entree{16592}{4:36:32}{Une petite sous-catégorie $C$ est génératrice par épimorphismes stricts universels si tout objet est but d'une famille universellement strictement épimorphique issue de $C$ ; c'est plus fort que strictement génératrice.}
\entree{16654}{4:37:34}{Dans un topos, toute famille épimorphique est strictement épimorphique et universelle.}
\entree{16654}{4:37:34}{Si $E$ vérifie $L$, a des produits fibrés, des noyaux $(u,v)\mapsto\mathrm{Ker}(u,v)$ formant un foncteur accessible, et une telle $C$, alors tout objet de $E$ est accessible et $E$ est accessible (SGA 4, I, 9.11).}
\entree{16782}{4:39:42}{Exemple : tout topos est une catégorie accessible.}
\entree{16782}{4:39:42}{Un ensemble $X$ de cardinal ${\ge}\,2$ n'est pas un objet accessible de $\mathbf{Ens}^{\mathrm{op}}$ ; donc $\mathbf{Ens}^{\mathrm{op}}$, qui a petites limites inductives et un générateur, n'est pas accessible.}
\entree{16908}{4:41:48}{Pour $\pi$ infini, les parties de cardinal ${\le}\,\pi$ d'un ensemble $I$ forment un ensemble ordonné dominant $\pi$, et $\prod_{i\in I}Y_i$ est la limite projective des produits partiels correspondants.}
\entree{17102}{4:45:02}{Si $\mathrm{card}\,I>\pi$ et $\mathrm{card}\,Y_i\ge 2$, l'indicatrice d'un point de $\prod Y_i$ ne se factorise par aucun produit partiel ; l'opposée d'un topos non dégénéré, ou de $\mathbf{Ab}$, n'est pas accessible.}
\entree{17297}{4:48:17}{Question : si $E$ et $E^{\mathrm{op}}$ sont accessibles, $E$ est-elle équivalente à la catégorie ponctuelle ?}
\entree{17429}{4:50:29}{Définition : pour $C$ petite, $\mathrm{Ind}(C)$ est la catégorie des systèmes inductifs filtrants de $C$, identifiée par $\varinjlim$ à une sous-catégorie pleine de $\widehat{C}$.}
\entree{17557}{4:52:37}{$\mathrm{Hom}_{\mathrm{Ind}(C)}((X_i),(Y_j))=\varprojlim_i\varinjlim_j\mathrm{Hom}(X_i,Y_j)$ (SGA 4, I, 8).}
\entree{17684}{4:54:44}{$C\to\mathrm{Ind}(C)$ est universel : pour $E$ à limites inductives filtrantes, la restriction à $C$ des foncteurs $\mathrm{Ind}(C)\to E$ commutant à ces limites est une équivalence avec $\mathrm{Hom}(C,E)$.}
\entree{17880}{4:58:00}{Si $C$ a des limites inductives finies, $\mathrm{Ind}(C)$ a de petites limites inductives et, pour $E$ qui en a, $F\mapsto F|_C$ identifie les foncteurs commutant à ces limites aux foncteurs $C\to E$ exacts à droite.}
\entree{18008}{5:00:08}{Définition : $\mathrm{Ind}_\pi(C)$ est la sous-catégorie pleine de $\mathrm{Ind}(C)$ des systèmes indexés par un $I$ dominant $\pi$ ; pour $C$ petite, elle est $\pi$-accessible.}
\entree{18331}{5:05:31}{Un objet de $\mathrm{Ind}_\pi(C)$ est $\pi$-accessible si et seulement s'il est facteur direct d'un objet de $C$ ; un facteur direct d'un objet $\pi$-accessible l'est encore.}
\entree{18396}{5:06:36}{Définition : $C$ est karoubienne si tout projecteur a un facteur direct image ; $\mathrm{Ind}(C)$ et $\mathrm{Ind}_\pi(C)$ ne changent pas si l'on remplace $C$ par son enveloppe karoubienne.}
\entree{18460}{5:07:40}{Si $C$ est karoubienne et $E=\mathrm{Ind}_\pi(C)$, alors $C\simeq E_\pi$ et $E\simeq\mathrm{Ind}_\pi(E_\pi)$.}
\entree{18524}{5:08:44}{Pour $\pi'\ge\pi$, $\mathrm{Fil}^{\pi'}=\mathrm{Ind}_{\pi,\pi'}(C)$ (systèmes indexés par un $I$ dominant $\pi$ avec $\mathrm{card}\,I\le\pi'$) filtre $\mathrm{Ind}_\pi(C)$ par des sous-catégories essentiellement petites, et $\mathrm{Fil}^{\pi}\simeq C$.}
\entree{18589}{5:09:49}{Les objets de $\mathrm{Fil}^{\pi'}$ sont $\pi'$-accessibles ; si $\pi'^{\pi}=\pi'$, tout objet $\pi'$-accessible de $\mathrm{Ind}_\pi(C)$ est facteur direct d'un objet de $\mathrm{Fil}^{\pi'}$.}
\entree{18652}{5:10:52}{\q{Filtration cardinale} : famille $(\mathrm{Fil}^{\pi'})_{\pi'\ge\pi}$ de sous-catégories essentiellement petites, stables par $\varinjlim$ d'indice $I$ dominant $\pi$, $\mathrm{card}\,I\le\pi'$, tout objet de $\mathrm{Fil}^{\pi''}$ étant limite d'objets de $\mathrm{Fil}^{\pi'}$ d'indice dominant $\pi'$, de cardinal ${\le}\,\pi''$.}
\entree{18847}{5:14:07}{Les $\mathrm{Ind}_{\pi,\pi'}(C)$ forment une filtration cardinale de $\mathrm{Ind}_\pi(C)$ si l'on se borne aux cardinaux tels que $\pi''^{\pi'}=\pi''$ (borne $\pi''^{\pi'}$ de SGA 4, I, 9).}
\entree{18911}{5:15:11}{Théorème : si $E$ a de petites limites inductives, sont équivalents : a) $E$ est $\pi$-accessible ; b) $E\simeq\mathrm{Ind}_\pi(C)$, $C$ petite ; c) $E_\pi$ est essentiellement petite et strictement génératrice.}
\entree{19371}{5:22:51}{Si $C\subset E_\pi$ est petite et pleine, le prolongement $\mathrm{Ind}_\pi(C)\to E$ est pleinement fidèle ; c'est une équivalence si $C$ est strictement génératrice et stable par limites à au plus $\pi$ flèches.}
\entree{19433}{5:23:53}{Une limite inductive indexée par une petite catégorie $I$ est la limite, suivant l'ensemble ordonné (dominant $\pi$) des sous-catégories de $I$ à au plus $\pi$ flèches, des limites partielles.}
\entree{19702}{5:28:22}{Si $E$ est $\pi$-accessible et a de petites limites inductives, $E\simeq\mathrm{Ind}_{\pi'}(E_{\pi'})$ pour tout $\pi'\ge\pi$.}
\entree{19912}{5:31:52}{Un foncteur $\pi$-accessible ne transforme pas nécessairement objets $\pi$-accessibles en objets $\pi$-accessibles (exemple : $E$ ponctuelle).}
\entree{19978}{5:32:58}{Si $E$, $E'$ sont $\pi_0$-accessibles à petites limites inductives, $F\colon E\to E'$ est accessible si et seulement s'il existe $\pi\ge\pi_0$ et $f\colon E_\pi\to E'_\pi$ avec $F\simeq\mathrm{Ind}_\pi(f)$.}
\section{11 juillet 1973}\label{s-1973-07-11}
\video{4gynb9uTEE4}
\noindent\textit{Vidéo 3/4 (\href{https://www.youtube.com/watch?v=4gynb9uTEE4}{\texttt{4gynb9uTEE4}}) : bande 1 (5:41:33--6:46:44); bande 2, face A (6:46:44--7:37:01).}\par\medskip
\face{Bande 1}
\entree{20626}{5:43:46}{Pour $i\colon C\to E$, $\varinjlim_{(Y,u)\in C/X}i(Y)\to X$ est un isomorphisme pour tout $X$ si et seulement si $E\to\widehat{C}$, $X\mapsto\mathrm{Hom}(i(-),X)$, est pleinement fidèle : $i$ est dit strictement générateur.}
\entree{20819}{5:46:59}{Alors tout $X$ est but d'une famille strictement épimorphique $i(Y_\alpha)\to X$ ; réciproque vraie si $i$ est pleinement fidèle, $E$ à produits fibrés et à familles strictement épimorphiques universelles.}
\entree{21008}{5:50:08}{Si $E$ a de petites limites inductives et $C$ est petite, $i$ se prolonge en $i_!\colon\widehat{C}\to E$, adjoint à gauche de $E\to\widehat{C}$ ; $i$ est strictement générateur si et seulement si $i_!$ fait de $E$ une catégorie de fractions de $\widehat{C}$.}
\entree{21139}{5:52:19}{Une catégorie à petites limites inductives a une petite sous-catégorie strictement génératrice si et seulement si elle est équivalente à une sous-catégorie pleine réflexive d'un $\widehat{C}$, $C$ petite.}
\entree{21201}{5:53:21}{Exemples : $\mathbf{Ens}^{\mathrm{op}}$ (algèbres de Boole complètes atomiques) et $\mathbf{Ab}^{\mathrm{op}}$ (groupes abéliens compacts) ont petites limites inductives et une petite famille strictement génératrice.}
\entree{21464}{5:57:44}{Dans la catégorie des espaces compacts, l'espace ponctuel n'est pas strictement générateur : le foncteur ensemble sous-jacent est fidèle mais pas pleinement fidèle.}
\entree{21923}{6:05:23}{Pour $\Delta_D$ (limites projectives de type $D'$ seules), l'enveloppe de $C$ est la sous-catégorie de $\mathrm{Hom}(C,\mathbf{Ens})^{\mathrm{op}}$ engendrée par $C$ par limites de type $D'$ itérées ; $C$ s'y plonge pleinement.}
\entree{22054}{6:07:34}{Si $R_0$ représente une théorie pour $\Delta_D$, la théorie prolongée à $\Delta$ est représentée par $R=T(\mathbf{Ens})^{\mathrm{op}}$, enveloppe de $R_0$ pour les petites limites projectives ; $R_0\subset R$ est pleine.}
\entree{22118}{6:08:38}{Une théorie relative à $\Delta_D$ équivaut au couple $(S,\Sigma_0)$, $S=T(\mathbf{Ens})$, $\Sigma_0=R_0^{\mathrm{op}}\subset S$, où $S$ est l'enveloppe de $\Sigma_0$ pour les petites limites inductives.}
\entree{22244}{6:10:44}{Pour $D'$ formé des diagrammes finis, les théories algébriques sont les petites catégories $R_0$ à limites projectives finies ; dualement, $\Sigma_0$ est petite à limites inductives finies.}
\entree{22374}{6:12:54}{Définition : $X$ est de présentation finie ($\aleph_0$-accessible) si $\mathrm{Hom}(X,-)$ commute aux limites inductives filtrantes ; on note $S_{\mathrm{pf}}$ la sous-catégorie pleine de ces objets.}
\entree{22438}{6:13:58}{Pour une théorie définie par limites projectives finies, $\Sigma_0\simeq S_{\mathrm{pf}}$ et $S\simeq\mathrm{Ind}(\Sigma_0)$ ; $S$ et $\Sigma_0$ se déterminent mutuellement, et $R_0\simeq(S_{\mathrm{pf}})^{\mathrm{op}}$.}
\entree{22501}{6:15:01}{Pour $\Sigma_0\subset S$ pleine, $\mathrm{Ind}(\Sigma_0)\to S$ est pleinement fidèle si et seulement si $\Sigma_0\subset S_{\mathrm{pf}}$ ; si $\Sigma_0$ est stable par limites inductives finies, c'est une équivalence si et seulement si $\Sigma_0$ est strictement génératrice.}
\entree{22570}{6:16:10}{Exemple : dans les catégories de groupes ou de modules, les objets de présentation finie sont les groupes ou modules de présentation finie au sens usuel.}
\entree{22638}{6:17:18}{$S$ à petites limites inductives est $T(\mathbf{Ens})$ pour une théorie définie par limites projectives finies si et seulement si $S_{\mathrm{pf}}$ est essentiellement petite et strictement génératrice.}
\entree{22766}{6:19:26}{Dans $\mathbf{Gr}$, pour $C$ pleine d'objet $\mathbb{Z}$, $\widehat{C}$ est formée des ensembles où opère le monoïde $(\mathbb{Z},\times)$ ; $G\mapsto(G,x\mapsto x^n)$ est fidèle non plein : $\mathbb{Z}$ est générateur, non strictement.}
\entree{22960}{6:22:40}{Pour $C\subset E$ pleine, $E\to\widehat{C}$ pleinement fidèle, fidèle, conservatif, ou conservatif sur les monomorphismes : notions de génération distinctes, qui coïncident dans un topos.}
\entree{23147}{6:25:47}{Les groupes libres de type fini forment une petite sous-catégorie strictement génératrice de $\mathbf{Gr}$.}
\entree{23211}{6:26:51}{Les morphismes $T\to T'$ de théories définies par limites projectives finies correspondent aux foncteurs $S'\to S$ commutant aux petites limites inductives et préservant les objets de présentation finie.}
\entree{23465}{6:31:05}{Pour les produits finis, l'enveloppe d'un objet $X$ est $(\mathbf{Ens}_{\mathrm{f}})^{\mathrm{op}}$, d'objets $X^n$ ; pour les produits de cardinal ${\le}\,\pi$, l'opposée des ensembles de cardinal ${\le}\,\pi$.}
\entree{23662}{6:34:22}{Une théorie à un objet de base définie par produits finis est une catégorie d'objets $\mathbb{N}$, $n$ étant puissance $n$-ième de $1$, donnée par les $H_n=\mathrm{Hom}(X^n,X)$ et leur substitution (\q{analyseurs} de Lazard).}
\entree{23791}{6:36:31}{Les réalisations d'une telle théorie (de Lawvere) dans une catégorie $E$ à produits finis sont les foncteurs vers $E$ commutant aux produits finis.}
\entree{23853}{6:37:33}{Pour un ensemble $I$ d'objets de base, on part de l'opposée de la catégorie des ensembles finis au-dessus de $I$ (objets $\prod_i X_i^{n_i}$, $n_i$ presque tous nuls).}
\entree{23917}{6:38:37}{Pour une théorie définie par produits finis, $\Sigma_0$ est stable par sommes finies et strictement génératrice dans $S$, mais cela ne caractérise pas les couples $(S,\Sigma_0)$.}
\entree{24044}{6:40:44}{$S$ à petites limites est la catégorie des modèles d'une théorie à produits finis d'objets de base $I$ si et seulement si ses relations d'équivalence sont effectives et $I$, formé d'objets de présentation finie, détecte les épimorphismes effectifs ($\mathrm{Hom}(i,f)$ surjectifs).}
\entree{24233}{6:43:53}{Alors $S\to\mathbf{Ens}^I$, $X\mapsto(\mathrm{Hom}(i,X))_{i\in I}$, a un adjoint à gauche et est monadique (\q{tripleable}).}
\face{Bande 2, face A}
\entree{24608}{6:50:08}{Théorème de Birkhoff : une sous-catégorie pleine d'une catégorie de structures définies par produits finis, stable par produits, sous-objets et quotients, est définie par des équations.}
\entree{24671}{6:51:11}{Une telle sous-catégorie est réflexive : l'adjoint associe à $X$ son plus grand quotient appartenant à la sous-catégorie.}
\entree{24737}{6:52:17}{Contre-exemple : les groupes abéliens sans torsion, stables par produits et sous-groupes mais non par quotients, ne sont pas définissables par équations dans $\mathbf{Ab}$.}
\entree{24866}{6:54:26}{À une sous-catégorie de Birkhoff $S'\subset S$ correspond $R\to R'$ bijectif sur les objets et surjectif sur les flèches : $T'$ se déduit de $T$ en égalisant des flèches.}
\entree{25511}{7:05:11}{Les groupes sont les monoïdes où une certaine flèche est inversible ; non stables par sous-monoïdes, ils ne sont pas définissables par équations dans les monoïdes.}
\entree{25649}{7:07:29}{Les sous-catégories pleines réflexives $S'\subset S$ à composé $S'\to\mathbf{Ens}^I$ monadique forment une classe plus large que celle des sous-catégories de Birkhoff.}
\entree{25846}{7:10:46}{Définition : un \q{pseudo-topos} vérifie les axiomes a), b), c) de Giraud : petites limites inductives et limites projectives finies ; sommes disjointes et universelles ; relations d'équivalence effectives et universelles.}
\entree{25975}{7:12:55}{Les foncteurs considérés commutent aux petites limites inductives et aux limites projectives finies : ce sont les images inverses $f^*$ des morphismes de topos de sens opposé.}
\entree{26048}{7:14:08}{Un topos $\mathcal{E}$ classifie une théorie $T$, d'objet universel $U\in T(\mathcal{E})$, si $f\mapsto f^*U$ est une équivalence $\mathrm{Top}(\mathcal{F},\mathcal{E})^{\mathrm{op}}\simeq T(\mathcal{F})$ pour tout $\mathcal{F}$ (2-flèches orientées par les images directes).}
\entree{26376}{7:19:36}{Pour $C$ petite à limites projectives finies, l'enveloppe de $C$ parmi les pseudo-topos est $\widehat{C}$, avec le plongement de Yoneda.}
\entree{26508}{7:21:48}{Définition : $f\colon C\to C'$ entre sites est continu si $F\mapsto F\circ f$ transforme faisceaux sur $C'$ en faisceaux sur $C$ ; cela ne dépend que de $\varepsilon_{C'}\circ f\colon C\to\widetilde{C'}$.}
\entree{26575}{7:22:55}{Par composition avec $\varepsilon_C$, les foncteurs $\widetilde{C}\to\widetilde{C'}$ commutant aux petites limites inductives s'identifient aux foncteurs continus $C\to\widetilde{C'}$.}
\entree{26704}{7:25:04}{Si $C$, $C'$ ont des limites projectives finies et $f$ commute à celles-ci, $f$ est continu si et seulement s'il transforme familles couvrantes en familles couvrantes.}
\entree{26704}{7:25:04}{Exemple : pour $X'\to X$ continue, l'image inverse des ouverts est un tel foncteur continu du site des ouverts de $X$ dans celui de $X'$.}
\entree{26767}{7:26:07}{En général, $f$ est continu si et seulement si, pour tout crible couvrant $R$ de $X$, $\varepsilon_{C'}f(X)=\varinjlim_{Y\in C/R}\varepsilon_{C'}f(Y)$.}
\entree{26832}{7:27:12}{Un crible $R$ de $X$ est couvrant si et seulement si $\varepsilon(X)=\varinjlim_{Y\in C/R}\varepsilon(Y)$ dans $\widetilde{C}$.}
\entree{27021}{7:30:21}{Si $C$, $C'$ ont des limites projectives finies et la topologie de $C'$ est moins fine que la canonique, $f$ continu induit un morphisme de topos si et seulement si $f$ est exact à gauche.}
\entree{27085}{7:31:25}{Définition : un morphisme de sites est un foncteur continu qui induit un morphisme de topos ; on l'oriente dans le sens de ce morphisme.}
\entree{27151}{7:32:31}{Pour un site $C$ à limites projectives finies et un topos $\mathcal{E}$, $\mathrm{Top}(\mathcal{E},\widetilde{C})^{\mathrm{op}}$ équivaut aux foncteurs $C\to\mathcal{E}$ exacts à gauche transformant familles couvrantes en familles épimorphiques.}
\entree{27279}{7:34:39}{$\widehat{C}$ est $\widetilde{C}$ pour la topologie la plus grossière (familles couvrantes : celles dont un membre a une section).}
\entree{27279}{7:34:39}{Si $C$ a des limites projectives finies, $\mathrm{Top}(\mathcal{E},\widehat{C})^{\mathrm{op}}$ équivaut aux foncteurs exacts à gauche $C\to\mathcal{E}$.}
\section{12 juillet 1973}\label{s-1973-07-12}
\video{ORG14LNb5gE}
\noindent\textit{Vidéo 4/4 (\href{https://www.youtube.com/watch?v=ORG14LNb5gE}{\texttt{ORG14LNb5gE}}) : bande 1 (0:00:00--0:49:57); bande 2, face A (0:49:57--1:38:22); bande 2, face B (1:38:22--2:26:53); bande 3, face A (2:26:53--3:15:55); bande 3, face B (3:15:55--4:05:06); bande 4, face A (4:05:06--4:53:49); bande 4, face B (4:53:49--5:36:07).}\par\medskip
\face{Bande 1}
\entree{5}{0:00:05}{Si $C$ est une petite catégorie à limites projectives finies, $\widehat{C}$ classifie la théorie $T_C$ : $\mathrm{Hom}_{\mathrm{Top}}(\mathcal{E},\widehat{C})\simeq\mathrm{Hom}_{\mathrm{lex}}(C,\mathcal{E})$ pour tout topos $\mathcal{E}$.}
\entree{131}{0:02:11}{Exemple : le topos classifiant des groupes est la catégorie des foncteurs covariants des groupes de présentation finie vers $\mathbf{Ens}$.}
\entree{195}{0:03:15}{Exemple : le topos classifiant des anneaux commutatifs est celui des foncteurs des anneaux de présentation finie vers $\mathbf{Ens}$, i.e. des préfaisceaux sur les schémas affines de type fini sur $\mathbb{Z}$.}
\entree{325}{0:05:25}{L'enveloppe $\widetilde{C}$ de $C$ par limites projectives finies est l'opposée de la catégorie des foncteurs de présentation finie $C\to\mathbf{Ens}$ ; $C\to\widetilde{C}$ est pleinement fidèle.}
\entree{388}{0:06:28}{Si $\mathcal{E}$ a des limites projectives finies, $\mathrm{Hom}(C,\mathcal{E})\simeq\mathrm{Hom}_{\mathrm{lex}}(\widetilde{C},\mathcal{E})$.}
\entree{454}{0:07:34}{Topos enveloppant d'une petite catégorie $C$ : le topos des préfaisceaux sur $\widetilde{C}$ ; les morphismes d'un topos $\mathcal{E}$ vers lui équivalent à $\mathrm{Hom}(C,\mathcal{E})$.}
\entree{603}{0:10:03}{Localisation : pour un ensemble $\Sigma$ de flèches d'un topos $\mathcal{E}$, on cherche un $\Delta$-foncteur universel $\mathcal{E}\to\mathcal{E}[\Sigma^{-1}]$ rendant inversibles les flèches de $\Sigma$.}
\entree{736}{0:12:16}{Les sous-topos de $\mathcal{E}$ correspondent bijectivement aux ensembles de flèches de $\mathcal{E}$ saturés (stables par limites et de nature locale) (SGA 4, exercice).}
\entree{801}{0:13:21}{Définition : un morphisme de topos $i\colon\mathcal{E}'\to\mathcal{E}$ est un plongement si $i_*$ est pleinement fidèle, i.e. si $i^*$ est un foncteur de localisation.}
\entree{930}{0:15:30}{Les sous-topos de $\mathcal{E}$ s'identifient aux sous-catégories strictement pleines dont l'inclusion admet un adjoint à gauche exact à gauche.}
\entree{997}{0:16:37}{Pour une petite catégorie $C$, les topologies $J$ sur $C$ correspondent bijectivement, via $J\mapsto\mathrm{Sh}(C,J)$, aux sous-topos de $\widehat{C}$.}
\entree{1059}{0:17:39}{Une application continue $f\colon X\to Y$ induit un plongement de topos ssi la topologie de $X$ est l'image réciproque de celle de $Y$ ; les sous-espaces de $Y$ donnent des sous-topos, mais pas tous.}
\entree{1124}{0:18:44}{Un morphisme $\mathcal{F}\to\mathcal{E}$ se factorise par un sous-topos $\mathcal{E}'$ ssi son image directe se factorise par $\mathcal{E}'$, ssi son image inverse se factorise par la localisation $\mathcal{E}\to\mathcal{E}'$.}
\entree{1189}{0:19:49}{À un sous-topos $\mathcal{E}'$ on associe les flèches rendues inversibles par $i^*$ ; à $\Sigma$, les objets $X$ tels que $\mathrm{Hom}(Z,X)\to\mathrm{Hom}(Y,X)$ soit bijectif pour tout $Y\to Z$ de $\Sigma$.}
\entree{1320}{0:22:00}{$\mathcal{E}[\Sigma^{-1}]$ existe et est un sous-topos de $\mathcal{E}$ ; tout sous-topos s'obtient ainsi en inversant une famille de flèches.}
\entree{1382}{0:23:02}{Si $\mathcal{E}=\mathrm{Sh}(C,J)$, ses sous-topos correspondent aux topologies sur $C$ plus fines que $J$ ; de même pour les sous-théories de la théorie représentée par $\mathcal{E}$.}
\entree{1967}{0:32:47}{Si une petite sous-catégorie pleine $C$ engendre le topos $\mathcal{E}$, $\mathcal{E}$ est un sous-topos de $\widehat{C}$ : $\mathcal{E}\simeq\mathrm{Sh}(C,J)$ pour la topologie induite.}
\entree{2248}{0:37:28}{Pour $C$ pleine dans $\mathcal{E}$ : (a) $\mathcal{E}\to\widehat{C}$ pleinement fidèle équivaut à (b) tout $X$ est limite inductive des $Y$ de $C/X$ ; (b) entraîne (c) : tout $X$ est but d'une famille strictement épimorphique issue de $C$.}
\entree{2380}{0:39:40}{Contre-exemple à (c) $\Rightarrow$ (b) : dans les groupes, avec $C=\{\mathbb{Z}\}$, $\mathbb{Z}*\mathbb{Z}$ vérifie (c) mais n'est pas limite inductive de $C/(\mathbb{Z}*\mathbb{Z})$.}
\entree{2446}{0:40:46}{Agrandir une famille strictement épimorphique peut lui faire perdre ce caractère (toutes les flèches $\mathbb{Z}\to\mathbb{Z}*\mathbb{Z}$ dans les groupes).}
\entree{2572}{0:42:52}{Si les produits fibrés existent et sont distributifs sur les sommes, $(X_i\to X)$ est strictement épimorphique ssi $\coprod X_i\to X$ l'est ; c'est faux en général.}
\entree{2765}{0:46:05}{Énoncé corrigé : (c$'$) tout $X$ est but d'une famille issue de $C$ strictement et universellement épimorphique ; alors (c$'$) entraîne (b).}
\entree{2829}{0:47:09}{Si une famille strictement et universellement épimorphique se factorise par une famille $(X_i\to Y)$, celle-ci est encore strictement et universellement épimorphique.}
\entree{2891}{0:48:11}{(b) équivaut à (b$'$) : pour tout $X$, la famille de toutes les flèches $Y\to X$, $Y$ dans $C$, est strictement épimorphique.}
\face{Bande 2, face A}
\entree{3001}{0:50:01}{Une famille $(Y\to X)$ est strictement épimorphique ssi pour tout $Z$, $\mathrm{Hom}(X,Z)$ s'identifie aux familles compatibles d'éléments des $\mathrm{Hom}(Y,Z)$ (condition de faisceau).}
\entree{3843}{1:04:03}{Théorème : $R\in\Delta$ représente une théorie algébrique ssi $R$ est un topos ; $B\colon I\to R$ est basique ssi les produits finis des $B(i)$ engendrent $R$.}
\entree{4018}{1:06:58}{Tout topos $R$ représente une théorie algébrique : si $C\subset R$ est petite, génératrice et stable par limites projectives finies, $T_R$ est une sous-théorie pleine de $T_C=\mathrm{Hom}_{\mathrm{lex}}(C,-)$.}
\entree{4236}{1:10:36}{Pour un tel $C$, les morphismes de topos $\mathcal{E}\to R$ correspondent aux foncteurs exacts à gauche $C\to\mathcal{E}$ envoyant familles couvrantes sur familles épimorphiques.}
\entree{4307}{1:11:47}{Si $R'\to R$ induit pour tout topos $\mathcal{E}$ un plongement pleinement fidèle $T_{R'}(\mathcal{E})\to T_R(\mathcal{E})$, $R'$ est un sous-topos de $R$ (exercice de SGA 4).}
\entree{4763}{1:19:23}{Un sous-topos de $R=\mathrm{Sh}(C,\tau)$ est $\mathrm{Sh}(C,\tau')$ avec $\tau'$ plus fine ; sa théorie se déduit de $T_R$ par un petit ensemble d'axiomes.}
\entree{5022}{1:23:42}{Les 2-limites projectives quelconques (non nécessairement filtrantes) de topos existent ; dans $\Delta$, ce sont des 2-limites inductives.}
\entree{5340}{1:29:00}{Construction : $C_i\subset R_i$ petites génératrices compatibles, $C$ la 2-limite inductive des $C_i$ dans $\mathbf{Cat}$, puis localisation du topos enveloppant de $C$ par un petit $\Sigma$.}
\entree{5405}{1:30:05}{La 2-limite inductive d'un système de petites catégories s'obtient en inversant les flèches cartésiennes de la catégorie (co)fibrée associée.}
\entree{5791}{1:36:31}{Une 2-limite projective de catégories accessibles, à foncteurs de transition accessibles, est accessible.}
\face{Bande 2, face B}
\entree{5906}{1:38:26}{La catégorie $J_0$ librement engendrée par produits finis sur un ensemble $I$ est l'opposée des familles d'ensembles finis indexées par $I$, presque tous vides ; c'est aussi l'enveloppe par limites projectives finies de $I$ discret.}
\entree{6037}{1:40:37}{$R_0=\widehat{J_0}$ classifie les familles d'objets indexées par $I$ : $\mathrm{Hom}_\Delta(R_0,\mathcal{E})\simeq\mathcal{E}^I$.}
\entree{6100}{1:41:40}{Ajout d'axiomes : $R_{\alpha+1}=R_\alpha[\Sigma_\alpha^{-1}]$ a les mêmes objets que $R_\alpha$, donc $J_0$ l'engendre encore.}
\entree{6232}{1:43:52}{Ajout de flèches données par un graphe $D=(H\rightrightarrows S)$ : $R_{\alpha+1}$ est la 2-somme amalgamée dans $\Delta$ de $R_\alpha$ et $\rho(\widetilde{D})$ sur $\rho(S)$ ($\rho$ topos enveloppant, $\widetilde{D}$ catégorie libre).}
\entree{6683}{1:51:23}{Dans cette construction, tout objet de $R_{\alpha+1}$ est couvert par des images d'objets de $R_\alpha$ ; donc $J_0$ engendre $R_{\alpha+1}$.}
\entree{7017}{1:56:57}{Étape limite : pour $\alpha$ limite, $R_\alpha$ est la 2-limite inductive dans $\Delta$ des $R_\beta$, $\beta<\alpha$ ; $J_0$ l'engendre encore.}
\entree{7147}{1:59:07}{Nécessité : si $I\to R$ est basique pour une théorie algébrique, $R$ est un topos engendré par les produits finis d'objets de $I$.}
\entree{7213}{2:00:13}{Suffisance : si les produits finis des objets de $I$ engendrent le topos $R$, la théorie s'obtient sans récurrence transfinie : produits finis, flèches entre eux, puis axiomes.}
\entree{7398}{2:03:18}{Pour $J$ petite à produits finis, la théorie $\mathcal{E}\mapsto\mathrm{Hom}_\Pi(J,\mathcal{E})$ des foncteurs préservant les produits finis est représentée par un topos $\rho_{\Pi,\Delta}(J)$.}
\entree{7732}{2:08:52}{Si la sous-catégorie $J_2\subset R$ des produits finis d'objets de $I$ engendre $R$, $R$ est un sous-topos de $\rho_{\Pi,\Delta}(J_2)$ ; deux localisations successives en font une.}
\entree{8191}{2:16:31}{La préservation des produits finis est un axiome : la flèche de constructions $\varepsilon_{X\times Y}\to\varepsilon_X\times\varepsilon_Y$ doit devenir un isomorphisme.}
\entree{8514}{2:21:54}{Si $R\in\Delta$ représente une théorie, sa catégorie des constructions est $R$ ; égaliser deux flèches revient à rendre inversible un noyau, d'où un seul type d'axiome.}
\face{Bande 3, face A}
\entree{8883}{2:28:03}{Dans un topos, $u\colon X\to Y$ est un épimorphisme ssi $Y\sqcup_X Y\to Y$ est un isomorphisme ; \q{être couvrant} s'écrit donc comme axiome d'inversibilité.}
\entree{9234}{2:33:54}{Le topos classifiant d'un objet sans structure est la catégorie des foncteurs des ensembles finis vers $\mathbf{Ens}$ (préfaisceaux sur $J=\{X^n\}$).}
\entree{9365}{2:36:05}{Si $B_T$ classifie la théorie $T$, sa catégorie de points équivaut à $T(\mathbf{Ens})$.}
\entree{9433}{2:37:13}{Pour tous topos $\mathcal{E}$ et $B$, $\mathrm{Hom}_\Delta(B,\mathcal{E})$ admet les limites inductives filtrantes, car dans un topos elles commutent aux limites projectives finies.}
\entree{9568}{2:39:28}{Les points essentiels du classifiant d'objets correspondent aux ensembles finis.}
\entree{9631}{2:40:31}{Le plus petit sous-topos contenant une famille de points d'un topos est défini par les flèches dont les fibres en ces points sont des isomorphismes.}
\entree{9758}{2:42:38}{Les sous-objets de l'objet final du classifiant d'objets $B$ sont $\varnothing$, l'image de $X\to e$ et $e$ : $B$ a trois sous-topos ouverts.}
\entree{9888}{2:44:48}{Le fermé complémentaire $\Phi_0$ de $B^*$ classifie les objets initiaux ; il est équivalent au topos ponctuel, de point l'ensemble vide.}
\entree{10017}{2:46:57}{L'ouvert $B^*$ classifie les objets $x$ tels que $x\to e$ soit épimorphique ; ses points sont les ensembles non vides.}
\entree{10145}{2:49:05}{Le sous-topos vide classifie la théorie qui n'a de modèle que dans le topos vide.}
\entree{10210}{2:50:10}{Le sous-topos $\Phi_{01}$ classifie les sous-objets de l'objet final ($x\to e$ monomorphisme) ; ses points sont les ensembles à au plus un élément.}
\entree{10342}{2:52:22}{$\Phi_1=\Phi_{01}\cap B^*$ classifie les objets finals ; $\Phi_0$ et $\Phi_1$ sont équivalents au topos ponctuel mais sont deux points distincts de $B$.}
\entree{10537}{2:55:37}{$\Phi_{01}\simeq\widehat{\Delta_1}$, où $\Delta_1=\{0\to 1\}$ : les morphismes $\mathcal{E}\to\widehat{\Delta_1}$ correspondent aux monomorphismes $x\hookrightarrow e$.}
\entree{10665}{2:57:45}{$\widehat{\Delta_1}$ est le topos de l'espace de Sierpiński (spectre d'un anneau de valuation discrète) ; un morphisme vers lui équivaut à un ouvert du topos source.}
\entree{10727}{2:58:47}{Tout espace topologique $T_0$ se plonge dans un produit de copies de l'espace de Sierpiński.}
\entree{10793}{2:59:53}{Les sous-topos d'un topos forment un treillis complet (intersections et bornes supérieures quelconques).}
\entree{11198}{3:06:38}{Dans l'opposée des groupes de présentation finie, le groupe universel est $\mathbb{Z}$ muni de sa structure de cogroupe ; le classifiant des groupes est son topos de préfaisceaux.}
\entree{11403}{3:10:03}{$\widehat{C}$ classifie les foncteurs $C\to\mathcal{E}$ dont le prolongement continu à $\widehat{C}$ est exact à gauche (foncteurs plats) ; si $C$ a des limites projectives finies, ce sont les foncteurs exacts à gauche.}
\entree{11533}{3:12:13}{Pour un groupe $G$ vu comme catégorie $C$ à un objet, $\widehat{C}=B_G$ ; un foncteur $C\to\mathcal{E}$ est un objet $X$ à opération de $G$, prolongé par $P\mapsto P\wedge^G X$.}
\entree{11601}{3:13:21}{Produit contracté : $P\wedge^G X$ est le quotient de $P\times X$ par l'opération diagonale de $G$, $(pg,x)\sim(p,gx)$.}
\entree{11667}{3:14:27}{Ce prolongement est exact à gauche ssi $X$ est un torseur sous le groupe constant $G_{\mathcal{E}}$ : $B_G$ classifie les $G$-torseurs.}
\face{Bande 3, face B}
\entree{11758}{3:15:58}{Problème : pour un groupe $G$ d'un topos $S$, représenter sur les topos $X$ au-dessus de $S$ le 2-foncteur des torseurs sous l'image inverse de $G$.}
\entree{12140}{3:22:20}{Les objets de $S$ munis d'une action à gauche de $G$ forment un topos $B_G$ : le foncteur d'oubli $B_G\to S$ est fidèle et commute à toutes les limites.}
\entree{12268}{3:24:28}{$p^*\colon S\to B_G$ munit un objet de l'action triviale ; un morphisme de groupes $G\to G'$ induit par restriction un morphisme essentiel $B_G\to B_{G'}$.}
\entree{12397}{3:26:37}{Torseur universel : $G$ agissant sur lui-même par translations à gauche (objet de $B_G$), avec l'action à droite de $\overline{G}=p^*G$ par translations à droite.}
\entree{12798}{3:33:18}{Universalité : pour $q\colon X\to S$, les $q^*G$-torseurs sur $X$ correspondent aux morphismes $X\to B_G$ au-dessus de $S$, via le produit contracté.}
\entree{13187}{3:39:47}{Sur un champ, la catégorie des objets localement isomorphes à une section $s$, avec leurs isomorphismes, équivaut à celle des torseurs sous $\underline{\mathrm{Aut}}(s)$.}
\entree{13251}{3:40:51}{Les $\mathcal{O}_X$-modules localement libres de rang $n$, avec leurs isomorphismes, équivalent aux torseurs sous $\mathrm{GL}_n(\mathcal{O}_X)$.}
\entree{13379}{3:42:59}{Le topos $B_{\mathrm{GL}_n}$ au-dessus du classifiant des anneaux classifie les topos annelés munis d'un module localement libre de rang $n$ (avec isomorphismes).}
\entree{13442}{3:44:02}{Les anneaux locaux sont classifiés par le sous-topos défini par la topologie de Zariski sur les schémas affines de type fini sur $\mathbb{Z}$.}
\entree{13506}{3:45:06}{Les anneaux strictement locaux sont classifiés par un sous-topos plus petit, défini par une topologie plus fine (la topologie étale).}
\entree{13632}{3:47:12}{Les objets localement constants ne relèvent pas d'une théorie représentable : ils ne sont pas stables par limites inductives (espaces totalement discontinus).}
\entree{13696}{3:48:16}{Les revêtements de rang $n$ de l'objet final ne forment pas une structure algébrique : dans $\mathbf{Ens}$, les ensembles à $n$ éléments ne sont pas stables par limites filtrantes.}
\entree{13826}{3:50:26}{Avec les seuls isomorphismes pour morphismes, les objets localement isomorphes au constant $I_n$ à $n$ éléments équivalent aux torseurs sous le groupe constant $\mathfrak{S}_n$.}
\entree{13891}{3:51:31}{Le topos $B_{\mathfrak{S}_n}$ des $\mathfrak{S}_n$-ensembles classifie les revêtements de rang $n$ de l'objet final, avec leurs isomorphismes.}
\entree{13957}{3:52:37}{La connexité du spectre d'un anneau s'écrit dans tout topos (tout idempotent est $0$ ou $1$, et $0\neq1$) et définit un sous-topos du classifiant des anneaux.}
\entree{14022}{3:53:42}{\q{Avoir exactement $n$ composantes connexes} ne définit pas une théorie : la propriété n'est pas stable par limites inductives filtrantes.}
\entree{14279}{3:57:59}{Un anneau localement à $n$ composantes connexes équivaut à un couple : $I$ localement constant de rang $n$, et un anneau $A_0$ 0-connexe sur $I$, dont $A$ est l'image directe.}
\entree{14409}{4:00:09}{Le foncteur $(I,A_0)\mapsto A$ est pleinement fidèle si l'on prend pour morphismes d'anneaux ceux qui induisent des isomorphismes sur les $\pi_0$.}
\entree{14601}{4:03:21}{Le classifiant des anneaux localement à $n$ composantes se construit au-dessus de $B_{\mathfrak{S}_n}$, en classifiant un anneau 0-connexe sur le revêtement universel $I_n$.}
\face{Bande 4, face A}
\entree{14978}{4:09:38}{Le classifiant d'un anneau au-dessus d'un objet $I$ d'un topos $B$ existe ; si $I$ est somme de $n$ copies de l'objet final, c'est $R^n\times B$, $R$ classifiant des anneaux.}
\entree{15238}{4:13:58}{Si $\mathcal{E}$ est un topos et $C$ a des limites projectives finies, $C\to\mathcal{E}$ est exact à gauche ssi son prolongement $\widehat{C}\to\mathcal{E}$ commutant aux limites inductives l'est.}
\entree{15302}{4:15:02}{Réciproque : si $\mathcal{E}$ a une petite sous-catégorie strictement génératrice $C$ stable par limites finies et que cette propriété vaut, $\mathcal{E}$ est un topos.}
\entree{15494}{4:18:14}{Le topos classifiant d'une théorie définie par limites projectives finies est connexe et acyclique.}
\entree{15751}{4:22:31}{Topos élémentaire : catégorie à limites projectives finies, cartésienne fermée, munie d'un classifiant des sous-objets.}
\entree{15944}{4:25:44}{La fidélité d'une opération de $G$ sur $X$ n'est pas stable par image inverse : ce n'est pas une structure algébrique au sens des topos.}
\entree{16521}{4:35:21}{Une théorie représentée par un topos non vide peut n'avoir aucun modèle ensembliste : il existe des topos sans points.}
\entree{16587}{4:36:27}{Exemple (Deligne) : boréliens de $[0,1]$ modulo la mesure nulle, recouvrements dénombrables à mesure nulle près ; le topos obtenu n'a aucun point.}
\entree{16853}{4:40:53}{Classifiant des corps : schémas affines de type fini sur $\mathbb{Z}$, familles couvrantes engendrées par $\{D(f)\to X,\ V(f)\to X\}$ et par $X_{\mathrm{red}}\to X$.}
\entree{16919}{4:41:59}{Pour $X$ quasi-compact et quasi-séparé, la topologie constructible (ouverts : réunions de parties constructibles) rend $X$ compact, avec les mêmes points.}
\entree{16981}{4:43:01}{Pour $X=\mathrm{Spec}\,A$, la limite inductive des anneaux de fonctions des partitions constructibles finies (structures réduites) a pour spectre $X$ constructible, de fibres les corps résiduels.}
\entree{17046}{4:44:06}{Le classifiant des corps s'identifie aux préfaisceaux sur ces schémas compacts (morphismes constructibles) changeant sommes finies en produits ; le corps universel est l'anneau des fonctions.}
\entree{17111}{4:45:11}{Anneaux absolument plats : anneaux dont tout module est plat ; ce sont les anneaux réduits de dimension de Krull nulle.}
\entree{17241}{4:47:21}{La platitude d'un module sur un anneau d'un topos est stable par image inverse par un morphisme de topos (Deligne).}
\entree{17307}{4:48:27}{Critère : $M$ est plat sur $A$ si pour tout $U$ et tout idéal de type fini $J\subset A_U$, $M_U\otimes J\to M_U$ est un monomorphisme.}
\entree{17371}{4:49:31}{Ce critère est une famille dénombrable d'axiomes : pour chaque $n$, sur $U_n=A^n$, l'idéal engendré par les $n$ projections universelles.}
\face{Bande 4, face B}
\entree{17700}{4:55:00}{Pour un schéma relatif de présentation finie sur un topos annelé, deux notions de platitude (Hakim) : morphisme de topos annelés plat, ou localement issu d'un morphisme plat de schémas.}
\entree{18019}{5:00:19}{Dans $\mathbf{Ens}$, un produit quelconque d'épimorphismes est un épimorphisme (analogue de la discontinuité totale) ; ce n'est pas vrai dans tout topos.}
\entree{18084}{5:01:24}{Le morphisme $\mathcal{E}\to\mathrm{pt}$ est essentiel ssi le foncteur faisceau constant commute aux produits quelconques (analogue de la connexité locale).}
\entree{18149}{5:02:29}{Les topos $\widehat{C}$ ont ces deux propriétés ; qu'ils soient les seuls est énoncé sans démonstration. \q{Hypertopos} : ces topos, avec les morphismes essentiels.}
\entree{18411}{5:06:51}{Une petite catégorie karoubienne $D$ est déterminée par $\widehat{D}$ : c'est la catégorie de ses points essentiels.}
\entree{18604}{5:10:04}{La catégorie des points de $\widehat{C}$, $C$ à limites projectives finies, est $\mathrm{Hom}_{\mathrm{lex}}(C,\mathbf{Ens})\simeq\mathrm{Ind}(C^{\mathrm{op}})$.}
\entree{18678}{5:11:18}{Pour les hypertopos, le seul candidat classifiant de $T_C$ est le topos des préfaisceaux sur la catégorie des points de $\widehat{C}$, non petite en général (\q{pseudo-hypertopos}).}
\entree{18936}{5:15:36}{Les conjectures de Weil sont démontrées par Deligne (1973), dans le cadre de la cohomologie étale fondée sur les univers.}
\entree{18997}{5:16:37}{L'axiome des univers équivaut à : tout cardinal est majoré par un cardinal fortement inaccessible ; l'existence d'un tel cardinal n'est pas démontrable dans ZFC.}
\entree{19061}{5:17:41}{\q{Galaxie} : ensemble satisfaisant les axiomes d'univers, sauf la stabilité par réunion d'une famille indexée par un de ses éléments, qui est affaiblie.}
\entree{19969}{5:32:49}{Les conjectures standard se formulent sur les cycles algébriques modulo équivalence numérique, sans cohomologie particulière ; elles comprennent l'égalité des équivalences homologique et numérique.}
\entree{20035}{5:33:55}{Les conjectures standard permettraient de construire les motifs à partir des seuls cycles algébriques, indépendamment de toute cohomologie.}

\end{document}
